%!TEX TS-program = pdflatex
% use-hpc -- Automated Software Engineering (Springer) submission draft.
% Compile: pdflatex ase-use-hpc; bibtex ase-use-hpc; pdflatex ase-use-hpc; pdflatex ase-use-hpc
% Needs: sn-jnl.cls, sn-mathphys-num.bst, ase-references.bib, figures/*.pdf (generated by analysis/make_figures.py),
%        data/tab_funnel.tex (generated by `use-hpc-eval selection-tex`).
% Every number in Sections 6-7 is derived from data/ or docs/ase/generated/ by the scripts in docs/ase/analysis/;
% docs/ase/analysis/check_claims.py verifies the quoted values against those files.
\documentclass[pdflatex,sn-mathphys-num]{sn-jnl}
\usepackage{url}
\usepackage{placeins}
\usepackage[T1]{fontenc}
\usepackage[scaled=0.92]{inconsolata}   % clean monospace for listings and \texttt
\usepackage{textcomp}
\usepackage{upquote}
\usepackage{graphicx}%
\usepackage{multirow}%
\usepackage{amsmath,amssymb,amsfonts}%
\usepackage{amsthm}%
\usepackage{xcolor}%
\usepackage{booktabs}%
\usepackage{algorithm}%
\usepackage{algorithmicx}%
\usepackage[noend]{algpseudocode}%
\usepackage{listings}%
\usepackage{tikz}
\usepackage{tabularx}
\usepackage{microtype}
\usepackage[utf8]{inputenc}
\usepackage[english]{babel}
\usepackage{enumitem}
\usepackage{pifont}
\usepackage{needspace}
\usepackage{hyperref}
\usetikzlibrary{arrows.meta,positioning,shapes.geometric,shapes.misc,fit,backgrounds,calc}
\hypersetup{colorlinks=true, linkcolor=black, urlcolor=black, citecolor=black}

\theoremstyle{thmstyletwo}%
\newtheorem{remark}{Remark}%
\newtheorem{property}{Property}%
\theoremstyle{thmstylethree}%
\newtheorem{definition}{Definition}%

\raggedbottom
\renewcommand{\topfraction}{0.92}
\renewcommand{\bottomfraction}{0.6}
\renewcommand{\textfraction}{0.08}
\renewcommand{\floatpagefraction}{0.85}
\setcounter{topnumber}{3}
\setcounter{bottomnumber}{2}
\setcounter{totalnumber}{5}
\setlength{\textfloatsep}{12pt plus 2pt minus 4pt}
\setlength{\floatsep}{10pt plus 2pt minus 2pt}

% ---------------- shorthands
\newcommand{\tool}{\texttt{use-hpc}}
\newcommand{\code}[1]{\texttt{#1}}
\newcommand{\RQ}[1]{\textbf{RQ#1}}
\newcommand{\V}[1]{\textsf{V#1}}
\newcommand{\I}[1]{\textsf{I#1}}
\newcommand{\cmark}{\ding{51}}
\newcommand{\xmark}{\ding{55}}
\newcommand{\pmark}{\ding{109}}

\definecolor{modelbg}{RGB}{235,244,255}
\definecolor{modelbar}{RGB}{60,110,180}
\definecolor{enginebg}{RGB}{236,250,238}
\definecolor{enginebar}{RGB}{60,145,80}
\definecolor{llmbg}{RGB}{248,238,252}
\definecolor{llmbar}{RGB}{140,80,160}

% ---------------- listings: light gray, Inconsolata, restrained syntax colours
\definecolor{codebg}{gray}{0.955}
\definecolor{coderule}{gray}{0.80}
\definecolor{codecomment}{gray}{0.45}
\definecolor{codekw}{RGB}{35,70,140}
\definecolor{codestr}{RGB}{150,60,40}
\definecolor{codelit}{RGB}{110,60,150}

\lstdefinelanguage{json}{
  morestring=[b]",
  morekeywords={true,false,null},
  moredelim=[s][\color{codestr}]{"}{"},
  sensitive=true}

\lstset{
  basicstyle=\ttfamily\fontsize{7}{8.4}\selectfont,
  backgroundcolor=\color{codebg},
  frame=single, rulecolor=\color{coderule}, framerule=0.4pt,
  framesep=5pt, xleftmargin=5pt, xrightmargin=5pt,
  aboveskip=8pt, belowskip=6pt,
  columns=fullflexible, keepspaces=true, upquote=true,
  breaklines=false, showstringspaces=false, numbers=none,
  captionpos=t,
  keywordstyle=\color{codekw}\bfseries,
  stringstyle=\color{codestr},
  commentstyle=\itshape\color{codecomment},
  identifierstyle=\color{black}
}
\lstdefinestyle{jsonstyle}{keywordstyle=\color{codelit}\bfseries}  % true/false/null


% \lstdefinelanguage{json}{morestring=[b]", sensitive=true}
% \lstset{frame=l, framerule=2pt, framesep=6pt, xleftmargin=9pt, framexleftmargin=0pt,
%         basicstyle=\ttfamily\fontsize{6.6}{7.7}\selectfont, breaklines=false, columns=fullflexible,
%         keepspaces=true, showstringspaces=false, numbers=none, upquote=true,
%         commentstyle=\itshape\color{black!55}, captionpos=t, aboveskip=6pt, belowskip=4pt,
%         backgroundcolor=\color{modelbg}, rulecolor=\color{modelbar}}
% \renewcommand{\lstlistingname}{Listing}

% run-in paragraph headings
\makeatletter
\renewcommand\paragraph{\@startsection{paragraph}{4}{\z@}%
  {6pt \@plus 2pt \@minus 1pt}{-0.6em}{\normalfont\normalsize\bfseries}}
\makeatother

\begin{document}

\title[\texttt{use-hpc}: agent-configured, checked Slurm pipelines]{\texttt{use-hpc}: Agent-Configured, Statically Checked Slurm Pipelines for Large-Scale GPU Experiments}

\author*[1]{\fnm{Majid} \sur{Babaei}}\email{majid.babaei@ufv.ca}
\affil*[1]{\orgdiv{School of Computing}, \orgname{University of the Fraser Valley}, \orgaddress{\state{BC}, \country{Canada}}}


\abstract{Empirical software engineering (SE) studies of open-weight large language models (LLMs) increasingly run thousands of GPU jobs on shared high-performance computing (HPC) clusters. Coding agents can write the Slurm scripts, but configuration errors that reach the cluster waste GPU time and can silently corrupt results. We present \tool{}, an open-source Claude Code skill in which the agent only proposes a declarative experiment graph, and fixed, tested programs check and execute it: a validator before submission, guards inside every job, a reduced smoke run of every stage, and an audit of every job and step. We evaluate this checking layer with 641 injected faults, a requeue experiment on Slurm and a 44{,}352-run pilot study, and the agent on ten artifacts selected from 4{,}814 SE publications. The static checks caught 438 of 439 malformed graphs, but no duplicate step names and none of 137 well-formed graphs with a wrong meaning. Setup gates stopped all 11 failing submissions, and Slurm cancelled 87 dependent jobs before they used any GPU time. GPU telemetry exposed two jobs that the scheduler reported as running but that left their GPUs almost idle for 2.7 hours. Resumption after a requeue lost no work, but a failed requeue is reported as success. With the skill, the agent configured all ten artifacts to pass the static checks, and four of six reached a passing setup gate after documented human interventions. The checking layer catches the faults it targets cheaply but not semantic errors; whether the configured experiments reproduce published results remains open.}

\keywords{{LLM agents, Reproducibility, High-performance computing, Slurm, Experiment management, Static checking}}

\maketitle

% =========================================================================
\section{Introduction}\label{sec:intro}

Empirical SE research increasingly evaluates LLM-based techniques by running them on benchmarks: hundreds of tasks, such as those of HumanEval+~\cite{liu2023evalplus} or ClassEval~\cite{du2024classeval}, crossed with several models and several random seeds. Because hosted model APIs can change without notice, guidelines for such studies recommend open-weight models that can be pinned and rerun~\cite{sallou2024breaking,baltes2025guidelines}. Open weights, however, shift the compute cost to the researcher. Eight conditions, seven models, 264 tasks and three seeds already make 44{,}352 runs. In academia, a study of this size runs on a shared high-performance computing (HPC) cluster, where a batch scheduler, typically Slurm~\cite{yoo2003slurm}, places jobs on GPU nodes.

Such studies fail in ways that a script run on a workstation does not reveal. A job that starts on missing inputs can write plausible but meaningless results; a stage that requests no GPU can compute on the CPU for its whole allocation; a parameter grid that reaches its time limit stops half-way; and the link from a number in a paper to the runs that produced it is easily lost. These are not hypothetical concerns: research code often fails to rerun as published~\cite{collberg2016repeatability,pimentel2019jupyter,trisovic2022research}, and SE artifacts are often incomplete~\cite{timperley2021artifact,liu2024artifacts}. Coding agents could take over much of this configuration work. General-purpose agents such as SWE-agent~\cite{yang2024sweagent} and OpenHands~\cite{wang2025openhands} resolve repository-level tasks and can write scheduler scripts, and several benchmarks evaluate agents that set up and reproduce research code~\cite{bogin2024super,siegel2024corebench,starace2025paperbench,baek2026artisan}. On a shared cluster, however, an agent's mistakes also cost other users: it may poll the scheduler in a tight loop, install packages on the shared login node, or report success when only the first job has ended.

This article studies one way to make an agent's work on a cluster \emph{checkable}. In \tool{}, the agent only \emph{proposes}: it writes a small, declarative \emph{experiment graph}. Fixed programs, which the agent is instructed to copy and configure rather than rewrite, then \emph{check and execute} the graph: a validator before submission, guards inside every job, a reduced \emph{smoke} graph that runs every stage at minimal size, Slurm's own dependency semantics, and an audit that reports every job and step. Constraining an LLM agent with deterministic enforcement is not a new idea. Runtime enforcement frameworks such as AgentSpec~\cite{wang2026agentspec} apply it to the agent's actions, and ExecutionAgent~\cite{bouzenia2025executionagent} uses execution feedback to guide agents that set up and run arbitrary repositories. What is specific to \tool{} is its target: multi-stage GPU experiments on a shared, scheduled cluster, where a configuration fault is cheap to detect before submission and expensive to discover afterwards.

We ask two separate questions, because the evidence for each differs. First, what does the fixed checking layer catch, what does it miss, and what does it cost? We answer this question with a fault-injection study on real agent-written graphs, a requeue experiment on Slurm, and measurements from a large pilot study whose graph the tool's author wrote. Second, can an agent equipped with the skill configure unseen research artifacts? We answer this question only up to the setup stage, with ten artifacts mined from recent SE venues.

\Needspace{8\baselineskip}
\paragraph{Contributions}
\begin{enumerate}[label=\textbf{C\arabic*}, leftmargin=2.2em, itemsep=1pt]
\item \emph{A propose-then-check design} for agent-configured cluster experiments: an experiment-graph model with parameter matrices, four Slurm dependency kinds and file contracts; six pre-submission checks (\V{1}--\V{6}); and four in-job invariants (\I{1}--\I{4}), each stated with the assumptions under which it holds (Section~\ref{sec:approach}).
\item \emph{\tool{}}, an open-source Claude Code skill whose Slurm pipeline template has about 1{,}550 lines of Python and Bash, together with a replication package whose scripts regenerate the article's figures and derived numbers from committed data.
\item \emph{An evaluation of the checking layer}: a fault-injection study with 641 mutants of ten agent-written graphs and a requeue experiment on Slurm, which together uncovered three defects of the tool, and measurements of fault containment, concurrency, GPU telemetry and overhead in a 44{,}352-run pilot study (Section~\ref{sec:resA}).
\item \emph{A feasibility study of the agent} on ten SE artifacts selected through a documented mining funnel, with a record of every fault, human intervention and deviation from the published artifacts (Section~\ref{sec:resB}).
\end{enumerate}

The remainder of the article is organized as follows. Section~\ref{sec:background} introduces the Slurm primitives that \tool{} builds on and the notion of an agent skill. Section~\ref{sec:approach} presents the propose-then-check design: the six-phase workflow, the experiment graph, the pre-submission checks \V{1}--\V{6}, the in-job invariants \I{1}--\I{4}, the smoke graph, and the audit and claim-to-job table.  Section~\ref{sec:design} states the research questions, the cluster setting, the subjects and the analyses. Sections~\ref{sec:resA} and~\ref{sec:resB} report the results. Section~\ref{sec:discussion} discusses what the checks are for, where gates belong and what smoke runs cover. Section~\ref{sec:related} positions \tool{} against workflow systems, reproducibility infrastructure and agents that set up research code, and Section~\ref{sec:threats} discusses threats to validity and the limits of our evidence. Section~\ref{sec:concl} concludes and outlines the experiments that would turn this feasibility evidence into evidence of benefit.

%The article does \emph{not} show that the agent's pipelines reproduce published numbers, that an agent with the skill outperforms one without it, or that the skill benefits novice users. Section~\ref{sec:threats} discusses these limits, and Section~\ref{sec:concl} outlines the experiments that would address them.

% =========================================================================
\section{Background}\label{sec:background}

\paragraph{Slurm} Slurm~\cite{yoo2003slurm} is a widely used workload manager on academic clusters. \tool{} builds on five of its primitives:
\begin{itemize}[leftmargin=1.4em, itemsep=1pt]
\item \emph{Batch jobs.} \code{sbatch} queues a script with requested resources (time, CPUs, memory, and generic resources such as GPUs); a fair-share scheduler decides when it starts. Login nodes are shared and meant for editing and submission only.
\item \emph{Job arrays.} \code{-{}-array=0-N\%K} queues $N{+}1$ tasks of one script, at most $K$ running at once; each task reads its index from \code{SLURM\_ARRAY\_TASK\_ID}.
\item \emph{Dependencies.} \code{-{}-dependency=afterok:J} starts a job only after job (or array) $J$ succeeded; \code{afternotok} only after it failed; \code{afterany} after it ended in any state; \code{aftercorr} starts array task $i$ after task $i$ of the upstream array succeeded. With \code{-{}-kill-on-invalid-dep=yes}, a job whose dependency can no longer be satisfied is cancelled before it starts.
\item \emph{Signals and requeue.} \code{-{}-signal=B:USR1@900} signals the batch shell 900\,s before the time limit, and \code{scontrol requeue} puts a job back in the queue under the same ID, so that its dependants keep waiting for it.
\item \emph{GPU partitioning.} NVIDIA's Multi-Instance GPU (MIG) splits an H100 into isolated slices, e.g.\ \code{h100\_1g.10gb} (one of seven compute slices) or \code{h100\_3g.40gb}~\cite{nvidia2023mig}.
\end{itemize}
Cluster operators add rules of conduct on top of these mechanisms. On the clusters of the Digital Research Alliance of Canada, for example, users must not run heavy work on login nodes or poll the scheduler in tight loops, should install Python packages from a curated \emph{wheelhouse} of prebuilt wheels, and should keep active data on scratch storage, which is periodically purged~\cite{alliance2026docs}.

\paragraph{Agent skills} An agent skill~\cite{anthropic2025skills} is a folder of instructions, reference documents and scripts that a coding agent such as Claude Code loads on demand. Until the skill is needed, only its name and a short description occupy the agent's context; when a request matches the description, the agent reads the instructions and, from there, whichever reference documents and scripts the current step requires. A skill can therefore carry both a method, which tells the agent what to do, and programs, which check whether it was done.

% =========================================================================
\Needspace{12\baselineskip}
\section{Approach}\label{sec:approach}

\subsection{Overview}\label{sec:overview}

Figure~\ref{fig:overview} shows the workflow. It takes two inputs: a document that describes experiments (a paper, report or design document) and the source code that implements them. The work proceeds in six phases. In phases 1, 3 and 6 the agent writes an artifact; in phases 2, 4 and 5 fixed programs act on what it wrote.

\begin{figure}
    \centering
    \includegraphics[width=1\linewidth]{figures/overview.drawio.pdf}
    \caption{Workflow of \tool{}. Octagons: inputs. Rounded dashed boxes: the agent writes. Square solid boxes: fixed programs check, execute and record. Diagnostics and failures send the agent back to the graph, never to the checks.}
    \label{fig:overview}
\end{figure}



% \begin{figure}[t]
% \centering
% \begin{tikzpicture}[font=\scriptsize,
%   box/.style={draw, align=center, minimum height=9mm, text width=26.5mm, inner sep=2pt},
%   llm/.style={box, rounded corners=5pt, fill=llmbg, draw=llmbar, dashed, thick},
%   prog/.style={box, fill=enginebg, draw=enginebar, thick},
%   data/.style={box, fill=modelbg, draw=modelbar, chamfered rectangle, chamfered rectangle xsep=2pt},
%   arr/.style={-{Stealth[length=2mm]}, thick}]
% \node[data] (doc) at (0.1,0.6) {document\\(claims, design)};
% \node[data] (src) at (0.1,-0.6) {source code\\(entry points)};
% \node[llm]  (inv)   at (3.45,0) {\textbf{1} [agent]\\inventory, source map};
% \node[prog] (scf)   at (6.85,1.3) {\textbf{2} [program]\\scaffold \code{hpc/}};
% \node[llm]  (graph) at (6.85,0) {\textbf{3} [agent]\\experiment graph $G$};
% \node[prog] (val)   at (10.25,0) {\textbf{4} [program]\\checks \V{1}--\V{6},\\dry run, smoke $\pi(G)$};
% \node[prog] (run)   at (10.25,-1.85) {\textbf{5} [program]\\Slurm jobs, guards \I{1}--\I{4},\\audit, collect};
% \node[llm]  (guide) at (6.85,-1.85) {\textbf{6} [agent]\\guide, claim-to-job\\table (reviewed\\by a human)};
% \draw[arr] (doc.east) -- (inv); \draw[arr] (src.east) -- (inv);
% \draw[arr] (inv) -- (graph); \draw[arr] (scf) -- (graph); \draw[arr] (graph) -- (val);
% \draw[arr] (val) -- (run); \draw[arr] (run) -- (guide);
% \draw[arr, dashed] (val.north) to[bend right=25] node[above, font=\tiny]{diagnostics} (graph.north east);
% \draw[arr, dashed] (run.north west) to[bend left=12] node[above, font=\tiny, pos=0.5, sloped]{failures} (graph.south east);
% \end{tikzpicture}
% \caption{Workflow of \tool{}. Octagons: inputs. Rounded dashed boxes: the agent writes. Square solid boxes: fixed programs check, execute and record. Diagnostics and failures send the agent back to the graph, never to the checks.}
% \label{fig:overview}
% \end{figure}

\begin{enumerate}[leftmargin=1.4em, itemsep=1pt]
\item \emph{Inventory.} The agent lists the document's research questions, subjects, factors, side studies, metrics and reported artefacts, and maps each item to the script that produces it. When the document and the code disagree, the agent must ask the human which one is right.
\item \emph{Scaffold.} A script copies the pipeline template into the \code{hpc/} directory of the repository.
\item \emph{Configure.} The agent writes the experiment graph as one JSON file, a requirements file and, where needed, small service hooks, for example one that starts a model server inside the job.
\item \emph{Verify.} A self-test checks the graph without running anything, a dry run prints every \code{sbatch} command, and a smoke run executes a reduced graph on GPU nodes.
\item \emph{Execute.} The full graph is submitted, guards in every job enforce the invariants, an audit reports progress, and a collector turns the outputs into tables and figures.
\item \emph{Document.} The agent writes a guide whose centrepiece is a table that maps each claim of the document to the steps and files that produce it. This table is meant for human review; no program checks it (Section~\ref{sec:coverage}).
\end{enumerate}
The central design decision is this \emph{division of labour}: the agent writes configuration and documentation, and every check is template code. A mistake in the graph should therefore surface as a diagnostic or a failed step, not as a wrong number in a results table. The division is a convention, not a sandbox: the skill's instructions tell the agent to copy and configure the template instead of writing its own pipeline, but nothing prevents the agent from editing the template.

\subsection{The experiment graph}\label{sec:graph}

\begin{definition}[Experiment graph]\label{def:graph}
An experiment graph is a tuple
$G=(\langle s_1,\dots,s_n\rangle, E, \mathit{cells}, \sigma, \rho)$, where
\begin{itemize}[leftmargin=1.4em, itemsep=1pt]
\item $\langle s_1,\dots,s_n\rangle$ is the list of \emph{stages} in file order;
\item $E\subseteq \{(s_i,s_j,k) \mid i<j,\ k\in\{\mathit{ok},\mathit{notok},\mathit{any},\mathit{corr}\}\}$ are the \emph{dependencies}, each from an earlier to a later stage, of one of Slurm's four kinds;
\item $\mathit{cells}(s)=\prod_{p\in P_s} D_p$ is the Cartesian product of the stage's \emph{matrix} parameters in a fixed order; a stage without a matrix has one empty cell, and a stage may reuse the cells of another stage (\code{matrix\_from});
\item $\sigma(s)$ is the ordered list of \emph{steps}, each a shell command template or a built-in action;
\item $\rho(s)$ are the stage's \emph{resources}: GPU kind and count (a whole GPU or a MIG slice), nodes, CPUs, memory, time limit and a requeue flag.
\end{itemize}
\end{definition}

Each stage becomes one \code{sbatch} submission, which is a job array when the stage has more than one cell. Each dependency becomes a Slurm dependency of the same kind on the upstream job ID, submitted with \code{-{}-kill-on-invalid-dep=yes}. The four kinds play distinct experimental roles. An \emph{ok} edge makes the upstream stage a \emph{gate}, typically a setup stage that checks data and models and runs a one-item pre-flight. An \emph{any} edge lets analysis and collection consume whatever data exist, so that one failed cell does not block the whole study. A \emph{corr} edge pipelines cell $i$ of one stage into cell $i$ of the next. A \emph{notok} edge runs a stage only when its upstream stage fails; it is supported but unused in our subjects. Command templates are filled from the graph's variables, the cell's parameter values, built-in values (such as the results directory, run ID and cell key) and the site settings in \code{cluster.env}.

\paragraph{Data contracts} A step $t$ may declare path patterns $N(t)$ (\emph{needs}) and $P(t)$ (\emph{produces}), relative to the run's results directory. Before $t$ runs, every pattern in $N(t)$ must match at least one file; otherwise $t$ fails with ``missing input''. After $t$ exits with code 0, every pattern in $P(t)$ must match; otherwise $t$ fails with ``declared output missing''. In the spirit of design by contract~\cite{meyer1992dbc}, a violated contract stops the step at which the fault becomes visible, instead of letting it surface several stages later. Because any single match satisfies a pattern, contracts check that outputs are present, not that they are complete. Contracts are optional, and Section~\ref{sec:rq5} reports how rarely the agent used them.

Listing~\ref{lst:graph} shows an excerpt of a graph that an agent session wrote for one of the artifacts of Section~\ref{sec:resB}. It contains a setup gate with critical checks, a two-cell retrieval stage, and a fine-tuning stage that consumes the retrieval outputs cell by cell through a \emph{corr} edge and a data contract.

\begin{lstlisting}[language=json, float=t, caption={Excerpt of the experiment graph an agent session wrote for RetriGen~\cite{zhang2025retrigen} (\code{studies/configs/isenglab\_\_retrigen/experiments.json}). Commands are shortened (\ldots); stage descriptions, the \code{collect} stage and the results section are omitted.}, label=lst:graph]
{"project": "r_retrigen", "require_gpu": true,
 "vars":  {"EPOCHS": "75", "WORK": "{RESULTS}/work"},
 "smoke": {"vars": {"EPOCHS": "1"}, "time": "00:55:00", "matrix_limit": 1},
 "stages": [
  {"name": "setup", "resources": {"gpus": "h100_1g.10gb:1", "time": "02:00:00"},
   "steps": [{"name": "check_gpu", "builtin": "check_gpu", "always": true},
             {"name": "model_present", "critical": true,
              "cmd": "... ollama list | grep -q 'codellama:7b-code-fp16'"},
             {"name": "codet5_snapshot", "cmd": "...", "critical": true,
              "produces": ["codet5-base/config.json"]}, ...]},
  {"name": "retrieve", "depends": [{"stage": "setup", "type": "afterok"}],
   "matrix": {"dataset": ["OldDataSet", "NewDataSet"]},
   "resources": {"gpus": "h100_3g.40gb:1", "time": "24:00:00", "requeue": true},
   "steps": [{"name": "prepare", "cmd": "...",
              "produces": ["work/{dataset}/dataset/{dataset}/*.csv"]}, ...]},
  {"name": "finetune", "depends": [{"stage": "retrieve", "type": "aftercorr"}],
   "matrix_from": "retrieve",
   "resources": {"gpus": "h100:1", "time": "72:00:00"},
   "steps": [{"name": "datas", "cmd": "python {HPC}/glue_retrigen.py jsonl {WORK}/{dataset} ...",
              "needs": ["work/{dataset}/hybrid/result/{dataset}/test/*alpha_0.5.csv",
                        "codet5-base/config.json"],
              "produces": ["work/{dataset}/datas/test.jsonl"]}, ...]}, ...]}
\end{lstlisting}

\subsection{Checks before submission}\label{sec:validate}

The self-test evaluates the graph on the login node, before any GPU is requested. Table~\ref{tab:checks} lists what it checks and, as the fault-injection study of Section~\ref{sec:rq1} established, what it does not. The checks are deliberately syntactic and structural: they establish that the graph is well formed and can be compiled into Slurm submissions, not that it encodes the intended experiment.

\begin{table}[t]
\centering\small
\caption{Pre-submission checks and in-job invariants of \tool{}: what is enforced, and what is assumed or not checked. ``Evidence'' points to the section that tests the item.}\label{tab:checks}
\begin{tabularx}{\textwidth}{@{}l >{\raggedright\arraybackslash}X >{\raggedright\arraybackslash}X l@{}}
\toprule
& Enforced by template code & Not checked / assumed & Evidence \\
\midrule
\V{1} Identity & stage names unique and file-safe & \emph{step} names unique within a stage (required by the schema documentation, but not enforced) & \S\ref{sec:rq1} \\
\V{2} Order & every dependency names an existing, earlier stage and a known kind; hence no cycles & whether the kind is the intended one & \S\ref{sec:rq1} \\
\V{3} Correspondence & \emph{corr} edges join stages with equal cell counts & -- & \S\ref{sec:rq1} \\
\V{4} GPU policy & every stage requests a GPU or slice unless marked CPU-only & whether the GPU kind, count or time limit is valid on the site & \S\ref{sec:rq1} \\
\V{5} Closure & every placeholder of every command and contract resolves, in every cell of the full and smoke graphs & whether it resolves to the intended value (graph and site variables share one namespace) & \S\ref{sec:rq1} \\
\V{6} Executability & every step has a command or built-in; an \code{sbatch} command builds for every stage & -- & \S\ref{sec:rq1} \\
\midrule
\I{1} No login-node work & the step runner exits (code 10) outside a Slurm job & -- & \S\ref{sec:rq1} \\
\I{2} GPU visible & the job exits (code 5) if no GPU is visible; GPU utilization is sampled every 60\,s & that the GPU is \emph{used}; MIG slices report no utilization & \S\ref{sec:rq1}, \S\ref{sec:rq3} \\
\I{3} Resumption & steps with a completion marker are skipped on restart & that commands skip records already written (append-and-skip); unique step names & \S\ref{sec:rq1} \\
\I{4} Requeue & on the pre-limit signal, the job stops its children and requeues itself, up to a cap & that a failed \code{scontrol requeue} is detected (it is not: the handler exits 0) & \S\ref{sec:rq1} \\
\bottomrule
\end{tabularx}
\end{table}

\subsection{Invariants inside jobs}\label{sec:invariants}

All stages share one batch script and one step runner (Algorithm~\ref{alg:runner}). The batch script refuses to run without a visible GPU (\I{2}), builds the Python environment from the wheelhouse, starts a GPU sampler that records utilization every 60\,s, runs optional service hooks, and finally starts the step runner in the background, so that the shell can handle the pre-limit signal immediately (\I{4}). The step runner, in turn, refuses to run outside a Slurm job (\I{1}).

\begin{algorithm}[t]
\caption{Step runner for cell $c$ of stage $s$ (inside a job).}\label{alg:runner}
\begin{algorithmic}[1]
\Require frozen graph $G$, results directory $R$, step ledger $L$
\If{not inside a Slurm job} \Return 10 \Comment{\I{1}}\EndIf
\State $\kappa \gets$ context(variables, $c$, built-ins);\quad $F\gets\emptyset$
\For{$t \in \sigma(s)$}
  \If{marker($s,c,\mathit{name}(t)$) exists and not $t.\mathit{always}$} \textbf{continue} \Comment{\I{3}}\EndIf
  \State write running-marker($s,c,t$) \Comment{read by the audit}
  \If{$\exists n\in N(t)$: glob($R$, fill($n,\kappa$)) $=\emptyset$} $rc\gets 6$ \Comment{missing input}
  \Else
    \For{attempt $=0..t.\mathit{retries}$} \State $rc\gets$ execute(fill($t,\kappa$)) \If{$rc=0$} \textbf{break} \EndIf \EndFor
    \If{$rc=0 \wedge \exists p\in P(t)$: glob($R$, fill($p,\kappa$)) $=\emptyset$} $rc\gets 7$ \Comment{output missing} \EndIf
  \EndIf
  \State append (job, task, restart, $s$, $c$, $t$, start, end, $rc$, note, command) to $L$
  \If{$rc=0$} write marker($s,c,\mathit{name}(t)$) \Else{} $F\gets F\cup\{t\}$; \textbf{if} $t.\mathit{critical}$ \textbf{then break} \EndIf
\EndFor
\State \Return $[F=\emptyset\,?\,0:1]$
\end{algorithmic}
\end{algorithm}

\begin{property}[Resumption]\label{prop:resume}
Assume that (a) step names are unique within each stage, (b) every step that is not marked \emph{always} is \emph{append-and-skip}, i.e.\ it writes one record per finished unit of work and skips units already recorded, and (c) a partially written last record is detected and ignored. Then, if a cell's job is interrupted and restarted any number of times, the set of units recorded after the final uninterrupted execution equals the set one uninterrupted execution records, and no step with a completion marker runs again.
\end{property}
\noindent The argument is short. A marker is written only after a step has returned 0 and its declared outputs exist, so a marked step is complete; an unmarked step reruns and, by (b), recomputes only the missing units. The property concerns the \emph{set} of units, not their values: units that were already recorded are not redrawn, so the values of stochastic units depend on where the interruption happened. Each assumption carries weight. Markers are keyed by step name, so without (a) a second step with the same name finds the first one's marker and is skipped, even on the first execution (Section~\ref{sec:rq1}). Assumption (b) is a property of the user's code, which the template cannot check; none of the ten artifacts of Section~\ref{sec:resB} was edited to establish it. One further condition is implicit: the requeue call must succeed. The released handler does not check that it does (Section~\ref{sec:rq1}).

\begin{remark}[Gates]\label{rem:gate}
If a stage fails, no stage that depends on it through \emph{ok} edges (or on the failed cell through \emph{corr} edges) starts: under \code{-{}-kill-on-invalid-dep=yes}, Slurm cancels such jobs, and cancellation propagates transitively. This is Slurm's semantics, not a property \tool{} establishes; \tool{}'s contribution is to emit the flag and the edges. Stages that depend through \emph{any} edges, such as collectors, do start (Section~\ref{sec:rq2}).
\end{remark}

\subsection{Smoke graph}\label{sec:smoke}

A configuration error in the last stage of a multi-day run should cost minutes, not days. The \emph{smoke graph} $\pi(G)$ keeps every stage, dependency and step of $G$, but overrides the size variables with small values (fewer tasks, seeds or items), keeps only the first $k$ cells of every matrix, and replaces every time limit with one short smoke limit. A passing smoke run has therefore executed every stage, step, dependency and declared contract at least once on real GPU nodes, with the real module stack, wheels, services and collector. It has not executed every \emph{cell}. In the pilot of Section~\ref{sec:resA}, the smoke graph reached one of seven models, and the only model-specific anomaly we observed lay outside it (Section~\ref{sec:rq3}). Smoke runs carry their own run-ID prefix, and their collected summary is labelled as not for reporting.

\subsection{Audit, collection and the claim-to-job table}\label{sec:coverage}

\paragraph{Audit} The audit joins two views of a run in one table. The scheduler's view comes from one \code{squeue} and one \code{sacct} call per snapshot; the pipeline's view comes from the step ledger, the running markers, the step logs and record counts. The table has one row per stage and cell, showing the state, node, run time against the time limit, steps completed, records written against records expected, and the current or failed step. Below the table, every failed step is listed with its exit code, reason and last log lines. The refresh interval cannot be set below 60\,s, and watch mode ends by itself once no job is pending or running, so monitoring cannot turn into a polling loop against the scheduler.

\paragraph{Collection} A final stage, attached with \emph{any} edges, runs the project's own analysis and then a built-in collector. The collector writes CSV tables; grouped summaries as CSV, Markdown and \LaTeX{}; figures; the step ledger; per-job GPU utilization; and a run manifest with the commit, content hashes of the code and configuration, and the frozen configuration. It then copies these files from purged scratch storage to a persistent location: by default the run's log directory inside the repository, or another directory set in \code{cluster.env}.

\paragraph{Claim-to-job table} In phase 6, the agent writes a guide around a table that maps each claim of the document (a number, table or figure, together with the question and metric it reports) to the steps and output files that produce it, and that lists the document's untested aspects as \emph{not covered}. This table is not a coverage guarantee. It is written by the agent and reviewed by a human, and no program checks that a claim's metric is computed from the files the table names. We therefore treat the table as documentation rather than as an evaluated contribution. %Sections~\ref{sec:impl-extend} and \ref{sec:concl} outline how it could be checked.

% % =========================================================================
% \section{Implementation}\label{sec:impl}

% This section describes how the design of Section~\ref{sec:approach} is realized in code. \tool{} is open source~\cite{usehpc2026repo} and has three layers: a \emph{package} that installs the skill into Claude Code, \emph{instructions} that tell the agent which method to follow, and a \emph{template} of programs that check and execute what the agent writes. We describe each layer in turn, following one run from installation to collected results, and close with the extensions that we consider most useful for SE researchers and developers. Table~\ref{tab:impl} maps every check and invariant of Section~\ref{sec:approach} to the code that realises it.

% The listings quote version 1.0.1 of the public repository (commit \code{2fc1509}), the version evaluated in Sections~\ref{sec:resA} and \ref{sec:resB}, abridged where marked (\ldots) and with some comments added.

% \begin{table}[t]
% \centering\footnotesize
% \caption{Where the elements of Section~\ref{sec:approach} are implemented (paths relative to the template \code{hpc/}; line counts at commit \code{2fc1509}).}\label{tab:impl}
% \setlength{\tabcolsep}{3pt}
% \begin{tabularx}{\textwidth}{@{}>{\raggedright\arraybackslash}p{2.2cm} >{\raggedright\arraybackslash}p{4.4cm} >{\raggedright\arraybackslash}X@{}}
% \toprule
% Element & Code & Mechanism \\
% \midrule
% \V{1}--\V{4}, \V{6} & \code{scripts/hpclib.py}: \code{validate} & one pass over the stages in file order (Listing~\ref{lst:validate}) \\
% \V{5} & \code{scripts/selftest.sh}, \code{hpclib.fill} & fills every command and contract of every cell of the full and the smoke graph; an unknown placeholder raises an error \\
% \V{6} (build) & \code{scripts/selftest.sh}, \code{slurm/submit.py} & a dry run of \code{submit.py} for the full and the smoke graph \\
% Smoke graph $\pi(G)$ & \code{hpclib.resolve} & overrides variables, truncates matrices to \code{matrix\_limit} cells, replaces time limits \\
% Edges, gates (Remark~\ref{rem:gate}) & \code{submit.py}: \code{sbatch\_cmd} & \code{-{}-dependency=}\textit{kind}\code{:}\textit{jobid} with \code{-{}-kill-on-invalid-dep=yes} (Listing~\ref{lst:submit}) \\
% \I{1} & \code{scripts/run\_stage.py} & exit code 10 when \code{SLURM\_JOB\_ID} is unset \\
% \I{2} & \code{slurm/common.sh}: \code{require\_gpu}, \code{start\_gpu\_monitor} & exit code 5 without a visible GPU; \code{nvidia-smi} sampled every \code{GPU\_MONITOR\_SECS} \\
% \I{3}, contracts & \code{scripts/run\_stage.py} & completion markers keyed by step name; glob checks of \emph{needs} and \emph{produces} (Listing~\ref{lst:runner}) \\
% \I{4} & \code{submit.py}, \code{common.sh}: \code{\_on\_timelimit} & \code{-{}-requeue -{}-signal=B:USR1@900}; the handler calls \code{scontrol requeue} (Listing~\ref{lst:job}) \\
% Audit & \code{scripts/audit.py} (317 lines) & one \code{sacct} and one \code{squeue} call per snapshot; refresh $\geq 60$\,s \\
% Collection, manifest & \code{scripts/collect\_results.py} (318), \code{scripts/manifest.py} (70) & tables, summaries, figures, pipeline records; commit and content hashes \\
% \bottomrule
% \end{tabularx}
% \end{table}

% \subsection{Packaging and installation}\label{sec:impl-pack}

% The repository is a Claude Code plugin \emph{marketplace} with a single plugin, \code{use-hpc}, which contains a single skill. Two JSON manifests declare the marketplace and the plugin; the skill itself lives in \code{plugins/use-hpc/skills/use-hpc/}. Listing~\ref{lst:use} shows the one-command installation and the commands the agent later runs inside a project. Where the plugin system is unavailable, an installer script copies the same skill folder to the user's personal skill directory (\code{\~{}/.claude/skills/use-hpc}). Nothing in the template depends on Claude Code: a user who copies \code{templates/hpc/} by hand obtains the same checks and the same job body.

% \begin{lstlisting}[language=bash, float=t, caption={Installing \tool{}, and the commands the agent runs in a project. Submissions other than dry runs require the user's consent.}, label=lst:use]
% claude plugin marketplace add ai4se4ai-lab/use-hpc-claude-skill
% claude plugin install use-hpc@ai4se4ai-hpc
% bash hpc/scripts/selftest.sh             # V1-V6 and dry runs on the login node; nothing submitted
% python3 hpc/slurm/submit.py --dry-run    # print every sbatch command
% python3 hpc/slurm/submit.py --smoke      # smoke graph: every stage at minimal size
% python3 hpc/slurm/submit.py              # the full graph
% bash hpc/status.sh --watch 120           # live audit, refresh >= 60 s
% \end{lstlisting}

% \subsection{Instructions: the method the agent follows}\label{sec:impl-instr}

% The skill folder holds an instruction file (\code{SKILL.md}, 104 lines), five reference documents (432 lines), the scaffold script and the template. The instruction file begins with a front matter whose \code{description} field decides when Claude Code loads the skill (Listing~\ref{lst:skill}). It lists the trigger situations explicitly, including requests that never mention a skill or HPC, such as ``run all the experiments'' of a shared paper. The body of the file states the deliverables, the ground rules of a shared cluster, and the six phases of Figure~\ref{fig:overview} as a to-do list, with a report to the user at the end of phases 1, 4 and 6.

% \begin{lstlisting}[caption={Excerpt of \code{SKILL.md}: the front matter that triggers the skill, the instruction to reuse the template, and two ground rules. Line breaks added.}, label=lst:skill]
% ---
% name: use-hpc
% description: Turn a project's source code plus a paper or design document that describes its
%   approach and experiments into a complete, GPU-only Slurm pipeline in the project's hpc/
%   directory. [...] Use this skill whenever the user wants to run, reproduce, scale out or test
%   experiments on an HPC cluster, [...]
% ---
% A tested, generic implementation ships with this skill in `templates/hpc/`. Copy it and
% configure it. Do not write a new pipeline from scratch: the template already handles
% requeueing, array tasks, contracts, the ledger, the GPU guard and the audit, and it was tested
% end to end.
% ## Ground rules (shared clusters)
% - **Login node = light work only**: editing, git, `sbatch`, reading logs, the audit, [...]
% - **Request only what is used.** Never poll Slurm faster than every 60 s. Never put `sleep` in
%   jobs. Ask before submitting anything big. A smoke run is small, but submitting it is still
%   the user's call.
% \end{lstlisting}

% The reference documents implement progressive disclosure: the instruction file names each one together with the situation in which to read it, so the agent loads only what the current task needs. \code{paper-to-experiments.md} carries the method of phases~1, 3 and~6 as four passes. Pass~1 builds the experiment inventory, a table of research questions with their decision rules, subjects, factors, design size, side studies, metrics, statistics, artefacts and fixed parameters. Pass~2 maps each inventory row to a source entry point, its output paths and its resumability. Pass~3 turns the inventory into a graph with a fixed shape: a setup gate that is the only stage with internet access, a main matrix with one cell per model and dataset (so that each job loads its model once), side studies that depend only on the setup gate, consumers attached by \emph{corr} or \emph{any} edges, and a final analysis-and-collection stage. Pass~4 produces the claim-to-job table of Section~\ref{sec:coverage}. The other references specify the graph schema (\code{experiments-schema.md}), Slurm and GPU usage, including multi-node jobs, requeueing and sizing (\code{slurm-gpu.md}), the audit and the collected outputs (\code{audit-and-results.md}), and site-specific notes for Alliance clusters (\code{alliance-cluster.md}).

% These instructions are the part of \tool{} that only the agent's compliance enforces. The template, described next, enforces the subset of them that a program can check: no computation outside a job, a GPU behind every stage, and no audit refresh faster than once per minute.

% \subsection{Scaffolding and site configuration}\label{sec:impl-config}

% In phase~2, \code{scaffold.sh} (23 lines) copies the 18 files of the template into \code{hpc/} and never overwrites an existing file unless the user passes \code{-{}-force}, so that rerunning it cannot destroy a configuration. Site-specific settings live in one Bash file, \code{config/cluster.env}: the Slurm account, the modules to load, the scratch directory for active work, the requeue settings (at most 20 requeues per job, with the warning signal 900\,s before the time limit), and the GPU sampling interval (60\,s). Every job sources this file, and the Python programs on the login node read it through \code{hpclib.source\_env}, which evaluates the file with Bash and keeps only the variables it defines. Both sides therefore see exactly the same values, including values computed from the environment, such as a default scratch path. At submission, \code{submit.py} freezes the site file, the graph, the requirements and the resolved graph into the run's log directory (\code{hpc/logs/<RUN\_ID>/config/}); every job and every later resubmission of the same run reads this frozen copy, so editing the configuration cannot change a study that is already running.

% \subsection{Graph resolution and static checks}\label{sec:impl-checks}

% All graph semantics live in one standard-library module, \code{hpclib.py} (229 lines). Its function \code{resolve} turns the graph that the agent wrote into an explicit one: it applies the smoke overrides when asked to, and it expands every matrix, or the matrix of the stage named by \code{matrix\_from}, into the list of its cells with \code{itertools.product}. In smoke mode it keeps the first \code{matrix\_limit} cells of every stage, replaces each stage's time limit by the smoke limit and, optionally, the array throttle; this is the whole implementation of $\pi(G)$.

% The function \code{validate} (Listing~\ref{lst:validate}) then makes a single pass over the stages in file order and returns a list of diagnostics; an empty list means that the graph is well formed. A set of the stages seen so far implements the order check \V{2}: a dependency on a stage that exists but has not yet been seen is reported as listed too late, which rules out cycles without a graph search. The listing also shows where the first defect of Section~\ref{sec:rq1} lies: \V{1} compares stage names, but no loop compares the names of the steps within a stage, although the schema documentation requires them to be unique.

% \begin{lstlisting}[language=Python, caption={\code{hpclib.validate}, abridged: the structural checks. Comments mark the checks of Table~\ref{tab:checks}.}, label=lst:validate]
% def validate(exp: dict) -> list[str]:
%     errs = []
%     names = [s.get("name") for s in exp.get("stages", [])]
%     if len(names) != len(set(names)):                                              # V1
%         errs.append("stage names must be unique")
%     seen = set()
%     for s in exp.get("stages", []):
%         n = s.get("name")
%         if not n or not re.fullmatch(r"[A-Za-z0-9_-]+", n):                         # V1
%             errs.append(f"stage name {n!r}: use letters, digits, _ and -")
%         for d in s.get("depends", []):
%             if d.get("stage") not in names:                                        # V2
%                 errs.append(f"{n}: depends on unknown stage {d.get('stage')!r}")
%             elif d["stage"] not in seen:                                           # V2
%                 errs.append(f"{n}: depends on {d['stage']!r}, which must be listed earlier ...")
%             if d.get("type", "afterok") not in ("afterok", "afterany", "afternotok", "aftercorr"):
%                 errs.append(f"{n}: dependency type {d.get('type')!r} unknown")      # V2
%             if d.get("type") == "aftercorr" and d.get("stage") in names:             # V3
%                 other = next(x for x in exp["stages"] if x["name"] == d["stage"])
%                 if "cells" in s and len(s["cells"]) != len(other.get("cells", [])):
%                     errs.append(f"{n}: aftercorr needs the same number of cells as ...")
%         r = s.get("resources", {})                                                 # V4
%         if exp.get("require_gpu", True) and not s.get("cpu_only") \
%                 and not (r.get("gpus") or r.get("gpus_per_node")):
%             errs.append(f"{n}: no GPU requested; require_gpu is on ...")
%         if not s.get("steps"):                                                     # V6
%             errs.append(f"{n}: no steps")
%         for st in s.get("steps", []):                                              # V6
%             if not st.get("name") or not (st.get("cmd") or st.get("builtin")):
%                 errs.append(f"{n}: every step needs a name and a cmd or builtin")
%         seen.add(n)
%     return errs
% \end{lstlisting}

% Placeholders are handled by two small functions (Listing~\ref{lst:fill}). The function \code{context} builds the namespace of one cell in four layers: the site variables from \code{cluster.env}, the graph's variables, the cell's matrix values, and the built-in values. A later layer overrides an earlier one. The function \code{fill} replaces every \code{\{NAME\}} with its value and raises an error for an unknown name; the regular expression skips \code{\$\{NAME\}}, which is left to Bash. Because a value may itself contain placeholders, such as \code{"WORK": "\{RESULTS\}/work"} in Listing~\ref{lst:graph}, \code{fill} repeats the substitution until nothing changes, up to five levels; this is the fix of the template defect T1 (Table~\ref{tab:faults}). The self-test (\code{selftest.sh}, 59 lines) implements \V{5} by resolving the full and the smoke graph and calling \code{fill} on the command and on every contract pattern of every step in every cell. It also checks the syntax of every shell script, compiles every Python file, and finishes with a dry run of \code{submit.py} for both graphs (\V{6}). The layering explains the second defect of Section~\ref{sec:rq1}: because the site variables form the bottom layer, a graph that loses its definition of \code{WORK} still resolves \code{\{WORK\}}, to the site's value, and \V{5} cannot notice the difference.

% \begin{lstlisting}[language=Python, caption={Placeholder resolution in \code{hpclib.py}, abridged. Site variables form the bottom layer of the namespace, which is why a removed graph variable can be silently shadowed by a site variable of the same name.}, label=lst:fill]
% _PH = re.compile(r"(?<![\$\\])\{([A-Za-z_][A-Za-z0-9_]*)\}")     # {NAME}, but not ${NAME}

% def context(exp, c, run_id, stage_name, cell, cell_idx):
%     ctx = {k: v for k, v in c.items() if re.fullmatch(r"[A-Za-z_][A-Za-z0-9_]*", k)}  # 1 site
%     ctx.update({k: str(v) for k, v in exp.get("vars", {}).items()})                  # 2 graph
%     ctx.update({k: str(v) for k, v in cell.items()})                                 # 3 cell
%     ctx.update({"RUN_ID": run_id, "RESULTS": ..., "STAGE": stage_name, ...})         # 4 built-in
%     return ctx

% def fill(text, ctx):
%     def rep(m):
%         k = m.group(1)
%         if k not in ctx:
%             raise KeyError(f"unknown placeholder {{{k}}} in: {text[:120]} ...")
%         return str(ctx[k])
%     out = str(text)
%     for _ in range(5):                       # values may contain placeholders (fix of T1)
%         new = _PH.sub(rep, out)
%         if new == out:
%             break
%         out = new
%     return out
% \end{lstlisting}

% \subsection{Compiling the graph to Slurm}\label{sec:impl-submit}

% The submitter, \code{submit.py} (171 lines), runs on the login node and does only light work. It resolves and validates the graph and refuses to submit (exit code 2) if \code{validate} reports anything or if the account in \code{cluster.env} is still the template's placeholder. It then freezes the configuration, writes the run manifest, and calls \code{sbatch} once per stage, in file order (Listing~\ref{lst:submit}). Because each stage is submitted after the stages it depends on, the job ID of every upstream stage is already known when a dependency is written; this is the operational reason for the order check \V{2}. The resource fields of a stage map one to one onto \code{sbatch} options, a stage with more than one cell becomes an array with an optional throttle, and a stage marked \code{requeue} receives \code{-{}-requeue} and the pre-limit signal. Every submission uses the same batch script, \code{stage.sh}, and passes the run ID and the stage name through the environment.

% \begin{lstlisting}[language=Python, caption={\code{submit.py}, abridged: one \code{sbatch} command per stage, and submission in file order.}, label=lst:submit]
% def sbatch_cmd(st, exp, c, run_id, deps, args):
%     r = st.get("resources", {}); n = len(st["cells"])
%     cmd = ["sbatch", "--parsable", f"--account={c.get('ACCOUNT', '')}",
%            f"--time={r.get('time', '01:00:00')}", ...]
%     if r.get("gpus_per_node"):
%         cmd.append(f"--gpus-per-node={r['gpus_per_node']}")
%     elif r.get("gpus"):
%         cmd.append(f"--gpus={r['gpus']}")
%     if r.get("requeue"):
%         cmd += ["--requeue", f"--signal=B:USR1@{c.get('REQUEUE_SIGNAL_SECS', '900')}"]
%     if n > 1:
%         thr = st.get("array_throttle")
%         cmd.append(f"--array=0-{n - 1}" + (f"%{thr}" if thr else ""))
%     if deps:
%         cmd += [f"--dependency={','.join(deps)}", "--kill-on-invalid-dep=yes"]
%     cmd.append(str(HPC / "slurm" / "stage.sh"))          # the one batch script of every stage
%     return cmd

% for st in exp["stages"]:                                  # in main(): file order
%     deps = []
%     for d in st.get("depends", []):
%         if d["stage"] in ids:                             # upstream stage already submitted
%             deps.append(f"{d.get('type', 'afterok')}:{ids[d['stage']]}")
%     out = subprocess.run(sbatch_cmd(st, exp, c, run_id, deps, a), capture_output=True, text=True)
%     ids[st["name"]] = out.stdout.strip().split(";")[0]
% \end{lstlisting}

% Two options support repair and inspection. With \code{-{}-run-id R -{}-only s1,s2}, the submitter reuses the frozen configuration of run \code{R} and resubmits only the named stages. Dependencies on stages that are not resubmitted are dropped with a note, and the data contracts take over their role: a resubmitted step whose inputs do not exist fails with ``missing input'' instead of computing on nothing. With \code{-{}-dry-run}, the submitter prints every \code{sbatch} command and submits nothing.

% \subsection{The job body: guards, environment and requeue}\label{sec:impl-job}

% Listing~\ref{lst:job} shows the batch script that every job runs and the handler of the pre-limit signal. The script first installs two traps: on exit, a line is appended to the run's job-ending ledger; on \code{USR1}, the requeue handler runs. The guard \code{require\_gpu} implements \I{2}: it exits with code 5 unless \code{CUDA\_VISIBLE\_DEVICES} is set and \code{nvidia-smi} lists a device, except in stages marked \code{cpu\_only}. The function \code{setup\_python} follows the Alliance practice: it loads the configured modules, creates a virtual environment in the job's local \code{\$SLURM\_TMPDIR}, installs the requirements from the wheelhouse with \code{pip install -{}-no-index}, and records \code{pip freeze} and the job's environment in the run's log directory. The GPU monitor runs \code{nvidia-smi} in sampling mode in the background and writes one CSV file per job and restart.

% \begin{lstlisting}[language=bash, caption={The job body (\code{slurm/stage.sh}, abridged) and the requeue handler (\code{slurm/common.sh}). The handler does not check the exit status of \code{scontrol requeue}.}, label=lst:job]
% source "$HPC_DIR/slurm/common.sh"     # stage.sh: the one batch script of every stage
% setup_traps                           # EXIT -> job-ending ledger;  USR1 -> _on_timelimit
% print_job_header
% require_gpu || exit 5                 # I2: no visible GPU, no computation
% setup_python || exit 4                # venv in $SLURM_TMPDIR, packages from the wheelhouse
% setup_distributed
% start_gpu_monitor                     # nvidia-smi every GPU_MONITOR_SECS
% run_hooks || exit 6                   # optional services, e.g. a private inference server
% run_stage                             # run_stage.py in the background, then wait

% _on_timelimit() {                     # common.sh: USR1, REQUEUE_SIGNAL_SECS before the limit (I4)
%     local n="${SLURM_RESTART_COUNT:-0}"
%     log_warn "time limit approaching (restart $n of max $MAX_REQUEUES)"
%     _cleanup                          # stop the runner, the GPU monitor and the services
%     if [ "$n" -lt "${MAX_REQUEUES:-20}" ]; then
%         _ending requeued
%         scontrol requeue "$REQUEUE_ID"
%         trap - EXIT; exit 0
%     fi
%     log_error "maximum number of requeues reached; work left unfinished"; exit 99
% }
% \end{lstlisting}

% Two details of the job body matter for correctness. First, the step runner is started in the background and the script waits for it, because Bash runs a trap only after the current foreground command has finished; a runner in the foreground would delay the requeue handler until the step ended, possibly beyond the time limit. Second, the handler requeues the job under its own ID (for an array task, \code{<array job>\_<task>}), so jobs that depend on it keep waiting instead of being released or cancelled. The last lines of Listing~\ref{lst:job} also show the third defect of Section~\ref{sec:rq1}: the exit status of \code{scontrol requeue} is ignored, so the job exits with code 0, and Slurm records it as completed, even when the requeue failed.

% Service hooks extend the job body without changing it. Before the steps run, the script sources \code{hooks/all.sh} and \code{hooks/<stage>.sh} if they exist. A hook may start a service in the background, such as an inference server for a local LLM, and registers its process ID so that the cleanup function stops it on exit and before a requeue. The template ships an example hook for Ollama that binds the server to a free port on the loopback interface, since other jobs may share the node, and waits for readiness with \code{curl -{}-retry} rather than a \code{sleep} loop. Three of the ten artifacts of Section~\ref{sec:resB} use Ollama to serve their models.

% \subsection{The step runner}\label{sec:impl-runner}

% The step runner, \code{run\_stage.py} (202 lines), executes the steps of one stage for one cell and implements Algorithm~\ref{alg:runner}. Listing~\ref{lst:runner} shows its core loop. Completion markers live on scratch storage, in \code{state/<stage>/<cell>/<step>.done}, and are named after the step, which is the root of the duplicate-name defect: two steps with the same name share one marker. Before a step runs, the runner writes a running marker that the audit reads to show the current step. Every execution of a step appends one row to the step ledger \code{steps.tsv} (time, job, array task, restart count, stage, cell, step, start, end, duration, exit code, note and the filled command), and the step's complete output goes to its own log file.

% \begin{lstlisting}[language=Python, caption={The core loop of \code{run\_stage.py}, abridged (Algorithm~\ref{alg:runner}).}, label=lst:runner]
% for s in st["steps"]:
%     marker = state / f"{s['name']}.done"                       # keyed by step name (I3)
%     if marker.exists() and not s.get("always"):
%         continue
%     running.write_text(json.dumps({"job": JOB, "task": TASK, "start": start, ...}))
%     try:
%         missing = [p for p in (H.fill(x, ctx) for x in s.get("needs", []))
%                    if not matches(results, p)]
%         if missing:
%             note, rc = "missing input: " + ", ".join(missing), 6
%         else:
%             for attempt in range(int(s.get("retries", 0)) + 1):
%                 if s.get("builtin"):
%                     rc = builtin(s["builtin"], ctx, log)         # check_gpu, check_nodes, collect
%                 else:
%                     rc = subprocess.run(["bash", "-o", "pipefail", "-c", cmd_text],
%                                         stdout=fh, stderr=subprocess.STDOUT,
%                                         env=env, cwd=str(H.REPO)).returncode
%                 if rc == 0:
%                     break
%             if rc == 0:
%                 absent = [p for p in (H.fill(x, ctx) for x in s.get("produces", []))
%                           if not matches(results, p)]
%                 if absent:
%                     note, rc = "declared output missing: " + ", ".join(absent), 7
%     except KeyError as e:                                       # unknown placeholder at run time
%         note, rc = f"config error: {e}", 8
%     ...                                                         # append one row to steps.tsv
%     if rc == 0:
%         marker.write_text(time.strftime("%F %T") + "\n")
%     else:
%         failed.append(s["name"])
%         if s.get("critical"):
%             break
% return 1 if failed else 0
% \end{lstlisting}

% The runner's exit codes identify the layer that stopped a step: 6 for a missing input, 7 for a missing declared output, 8 for an unknown placeholder at run time, 9 for any other exception in the runner, 10 for an attempt to run outside a Slurm job (\I{1}), and 11 for a cell index outside the stage's matrix. Commands run from the repository root under Bash's \code{pipefail} option, so a failure anywhere in a pipeline fails the step; in particular, a pipeline whose reader exits early fails with the writer's \code{SIGPIPE} status, 141, which is the exit code of one of the agent's configuration faults in Table~\ref{tab:faults}. Besides commands, a step may name one of three built-in actions: \code{check\_gpu} records the visible GPUs in the results directory, \code{check\_nodes} verifies in a multi-node job that every node has a GPU, sees the shared results directory and can resolve the master node, and \code{collect} runs the collector.

% \subsection{Audit, collection and manifest}\label{sec:impl-audit}

% The audit (\code{audit.py}, 317 lines) assembles its table from files and two scheduler calls (Listing~\ref{lst:audit}). It reads the job IDs that \code{submit.py} recorded in \code{jobs.tsv}, queries \code{sacct} and \code{squeue} once each for all of them, and combines each array task's scheduler state with the step ledger and the running markers of its cell. Where the two views disagree, the pipeline's view refines the scheduler's: a job that Slurm reports as \code{FAILED} is shown as \code{PARTIAL} when some of its steps succeeded, and a step that is still running is shown even when the scheduler has not yet reported the job. Progress bars come from a step's optional \code{progress} specification, which names an output file and an arithmetic expression for the expected number of records; the audit counts the file's JSON lines or CSV rows. Each snapshot is also saved without colour as \code{audit.txt} and \code{audit.json}, so the state of a run at any moment can be archived.

% \begin{lstlisting}[language=Python, caption={The audit's interaction with Slurm (\code{audit.py}, abridged): one \code{sacct} and one \code{squeue} call per snapshot, a lower bound on the refresh interval, and a watch mode that ends by itself.}, label=lst:audit]
% def slurm(ids):
%     for line in sh(["sacct", "-n", "-P", "-X", "-j", ",".join(ids),
%                     "--format=JobID,State,Elapsed,Timelimit,NodeList,ExitCode"]).splitlines(): ...
%     for line in sh(["squeue", "-h", "-r", "-j", ",".join(ids),
%                     "-o", "%i|%T|%r|%N|%M|%l"]).splitlines(): ...

% period = max(a.watch, 60) if a.watch else 0      # in main(): at most one refresh per minute
% while True:
%     data = assemble(run_id)
%     print(render(data, a.only, a.steps))
%     active = any(r["state"] in ("RUNNING", "PENDING", "REQUEUED") for r in data["rows"])
%     if not period or not active:
%         return 0                                 # leave watch mode when nothing is left to watch
%     time.sleep(period)                           # login-side viewer only, never inside a job
% \end{lstlisting}

% The collector (\code{collect\_results.py}, 318 lines) is driven by the \code{results} block of the graph. It concatenates matching CSV, JSON and JSON-lines files into one table each, computes grouped summaries (mean, median, sum, count, minimum, maximum, standard deviation or rate), renders them as CSV, Markdown and \LaTeX{} tables, and draws figures when matplotlib is available. It also exports the pipeline's own records: every step attempt, the latest attempt of every step, the submitted jobs, the job endings, per-job GPU statistics, and a completeness table taken from the audit. Directories with many small files can be packed into compressed tar archives, following the Alliance guidance for networked file systems. Finally, \code{manifest.py} (70 lines), which \code{submit.py} calls at every submission, records the Git commit and branch, the list of uncommitted files and a hash of the uncommitted diff, content hashes of the source directories and of \code{hpc/config}, the site variables except those whose names suggest secrets, and the times of all submissions to the run. Two result sets are comparable when these hashes match.

% \subsection{Extending the skill}\label{sec:impl-extend}

% The evaluation exposes concrete gaps, and the SE studies that motivate \tool{} raise needs that the current template does not meet. We describe the extensions we consider most useful, from the most direct to the most speculative. None of them is implemented or evaluated in this article.

% \paragraph{Closing the defects found here} Each of the three defects has a local fix. \code{validate} can compare step names within each stage, as it already does for stage names. \code{context} can record which layer supplied each value, so that the self-test warns when a placeholder used in the graph resolves only from \code{cluster.env}. The requeue handler can test the exit status of \code{scontrol requeue} and, if the requeue failed, record the failure in the job-ending ledger and exit with a non-zero code, so that dependants attached by \emph{ok} edges are cancelled instead of started on incomplete data. Each fix should come with the fault-injection or Slurm test that exposed the defect, run as a regression test against every release.

% \paragraph{Moving checks to where they are cheapest} Section~\ref{sec:rq2} shows that most gate failures were package-resolution problems that a login-node check finds in seconds. The self-test could resolve \code{requirements.txt} against the wheelhouse without installing anything (\code{pip install -{}-dry-run -{}-no-index}), as the evaluation harness already does. A second cheap check would link the data contracts statically: every \emph{needs} pattern of a step should be matched by a \emph{produces} pattern of an earlier step of the same stage or of an upstream stage, and a stage that reads another stage's outputs through an \emph{any} edge deserves a warning. Such a linkage check would turn some of the semantic faults that escaped all checks in Section~\ref{sec:rq1}, such as a removed dependency, into diagnostics, provided that the agent declares contracts, which it rarely did (Section~\ref{sec:rq5}). The skill's instructions could therefore require contracts on every step that crosses a stage boundary, and the self-test could report steps that lack them.

% \paragraph{Turning conventions into enforcement} Several guarantees rest on the agent following its instructions, most importantly that it does not edit the template. The scaffold script could record content hashes of the template files, and the self-test and the run manifest could compare against them, so that every edit becomes visible in the record of the run. Claude Code's hooks, which run before a tool call, could reject \code{sbatch}, \code{scancel} or polling commands that bypass \code{submit.py} and the audit, turning the ground rules of Listing~\ref{lst:skill} from instructions into policy.

% \paragraph{Better smoke coverage and telemetry} The first-$k$ projection reached only one of seven models in the pilot (Section~\ref{sec:smoke}). A projection that covers every value of every matrix parameter at least once, or every pair of values, would cost a few more cells and would have included the model whose jobs idled. The GPU sampler could compare utilization against a threshold and flag, or stop, a job that stays below it for a configurable time, so that idle GPUs become audit failures instead of a fact that only a later analysis finds.

% \paragraph{Other sites and runtimes} The scheduler-specific code is concentrated in a few places: \code{sbatch\_cmd} in the submitter, the \code{slurm} function of the audit, the requeue and restart logic in \code{common.sh}, and the \code{srun} launcher for multi-node steps. A small backend interface around these places would allow PBS Pro, LSF or Flux backends while keeping the graph, the checks and the runner unchanged. Four of the repositories excluded during manual screening needed Docker, which the cluster does not provide (Section~\ref{sec:subjects}); a step attribute that runs the command inside an Apptainer image~\cite{kurtzer2017singularity}, the container runtime that HPC sites commonly support, would admit such artifacts and would also pin their environments more tightly than a requirements file.

% \paragraph{Needs of SE researchers} Three needs recur in the SE studies we examined. First, LLM-based studies need model provenance: the manifest could record model digests or revisions and decoding parameters, and the collector could aggregate token counts and latencies as standard outputs, as SE guidelines for LLM studies ask~\cite{sallou2024breaking,baltes2025guidelines}. Second, campaigns such as mutation testing, fuzzing or test generation consist of thousands of tasks that each take seconds to minutes; one array task per cell is wasteful for them, and the Alliance asks users to bundle tasks shorter than about 15 minutes~\cite{alliance2026docs}, so an option that packs several cells into one task would be needed. Third, the claim-to-job table could become machine-checkable: a \code{claims} block in the graph could name, for each claim of the paper, the step and the collected file that produce it, and the self-test could check that each named step exists and declares a matching \emph{produces} pattern. Together with an export of the manifest and ledgers to a standard provenance format, this would give artifact evaluation committees a direct way to trace a reported number to the jobs that produced it.

% =========================================================================
\section{Evaluation design}\label{sec:design}

\subsection{Research questions}\label{sec:rqs}

The evaluation has two parts with different subjects. Separating them keeps results about the fixed template from being read as results about the agent.

\paragraph{Part A: the checking layer} The subjects are the ten agent-written graphs of Part B, which serve as seeds for fault injection, and a pilot study whose graph the tool's author wrote.
\begin{description}[leftmargin=1.1em, style=nextline, itemsep=2pt]
\item[\RQ{1} Static checking] Which configuration faults do the pre-submission checks detect before anything is submitted, which escape, and how does their cost grow with graph size?
\item[\RQ{2} Fault containment] For faults that reach the cluster, which layer reports them, how long after submission, and at what GPU cost?
\item[\RQ{3} Concurrency and observability] How close do the submitted graphs come to the concurrency their dependency structure permits, and what does per-job GPU telemetry show that scheduler states do not?
\item[\RQ{4} Overhead] What do the template's machinery and a smoke run cost?
\end{description}

\paragraph{Part B: the agent} The subject is a new Claude Code session with \tool{} for each artifact.
\begin{description}[leftmargin=1.1em, style=nextline, itemsep=2pt]
\item[\RQ{5} Feasibility on unseen artifacts] For ten artifacts of recent SE papers, can the agent produce graphs that pass the static checks and the setup gates; what configuration, glue code and human intervention did that take; and which faults surfaced, where?
\end{description}
Three questions are outside the scope of this article and are discussed in Section~\ref{sec:threats}: whether reruns reproduce the published numbers, how the skill compares with an agent without it or with a workflow system, and how much results vary across agent sessions.

\subsection{Setting}\label{sec:setting}

All runs used Nibi, a cluster of the Digital Research Alliance of Canada that is scheduled by Slurm 25.11. Its GPU nodes have up to four H100 80\,GB GPUs, many of them partitioned into MIG slices. We used one ordinary (default) allocation that competed with other users, so queue waits are part of the measurements. We express GPU cost in \emph{H100-equivalent hours}: the run time multiplied by the allocated compute share, where a whole H100 counts as 1, a \code{3g.40gb} slice as $3/7$ and a \code{1g.10gb} slice as $1/7$.

\subsection{Subjects}\label{sec:subjects}

\paragraph{Pilot (Part A)} The pilot is the evaluation of an approach for checking LLM-assembled agent teams, an unpublished study by the author of this article.
% TODO(author): cite the pilot study (AutoM2M) if a citable version exists; otherwise keep this description.
Its main matrix has 8 conditions $\times$ 7 open-weight code models (7B--32B parameters) $\times$ 264 tasks (HumanEval+~\cite{liu2023evalplus} and ClassEval~\cite{du2024classeval}) $\times$ 3 seeds $=$ 44{,}352 runs, plus side studies. Its graph has 26 jobs: a setup gate; 14 matrix jobs, one per model and benchmark; 6 side-study jobs; 3 repair jobs; 1 coding job; and 1 analysis job. All jobs except the analysis job, which the prototype ran without a GPU, use one or two H100s, and each job serves its model through a private inference server. The graph was written by the template's author, and the pipeline is the prototype from which the template was distilled: it has the same graph semantics, guards, ledgers and audit as the template, but stage-specific batch scripts instead of one shared script. The pilot therefore tests the design and its realisation on Slurm, not the agent. Its smoke graph has 10 jobs.

\paragraph{Artifacts (Part B)}
We selected six established software engineering venues to construct a focused, literature-grounded sample of recent research: the International Conference on Software Engineering (ICSE), the ACM International Symposium on the Foundations of Software Engineering (FSE), the International Conference on Automated Software Engineering (ASE), IEEE Transactions on Software Engineering (TSE), Empirical Software Engineering (EMSE), and ACM Transactions on Software Engineering and Methodology (TOSEM). Prior empirical surveys and bibliometric studies use overlapping sets of these venues as sources of high-quality software engineering research~\cite{silva2021topic,kapitsaki2026does}. Together, they cover general, automated and empirical software engineering research. We use them as a focused sampling frame, not as an exhaustive coverage of the field.

We retrieved all 4{,}814 publications of these venues from 2024 to 2026 from OpenAlex~\cite{priem2022openalex} and applied the scripted criteria of Table~\ref{tab:funnel}. A publication remained if it is a research paper on LLMs that links a GitHub artifact in its full text, and if that repository is reachable, uses open-weight models, and is a Python project with declared dependencies and documented run commands. The 118 remaining repositories were ranked by a fixed score that favours local model serving (a 10-point bonus for Ollama use), mentions of open-weight models, documented commands, recent activity and a licence. One researcher then screened repositories in rank order against manual criteria: the artifact must be the paper's own; quantitative results using open-weight models must be reported; a documented entry point must be available without Docker, which the cluster does not support; data must be public; models must have at most 34B parameters; the estimated experimental subset must require at most 300 H100-hours; and the study must not require human participants or paid services. Screening stopped when ten repositories passed. This required 24 screening decisions, including one inclusion later reversed after cloning revealed a call to a paid API. All criteria are defined in a single configuration file, and every decision is logged.


For each artifact, a new Claude Code session with \tool{} read the paper and the repository, cloned at the screened commit, and wrote the configuration (Table~\ref{tab:artifacts}). The human chose the executed subset and the comparison target of each artifact from the agent's inventory. For six artifacts, the configuration is a \emph{subset rerun}. For the other four, the session had reported defects in the documented commands that would block a rerun; their configuration is a \emph{probe} that runs every documented entry point once at minimal size, so that an exit code can confirm or refute each suspected defect.

\begin{table}[t]
\centering\small
\caption{Selection funnel of the artifact study (generated from \code{data/funnel.csv}).}\label{tab:funnel}
\input{data/tab_funnel.tex}
\end{table}

\subsection{Data and analyses}\label{sec:data-sources}

All numbers come from four sources: Slurm accounting (one \code{sacct} dump of the 188 pilot jobs of 11 submissions, and one of the 75 study jobs); the step ledgers and 60-s GPU samples of the pilot, summarised in committed CSV files (the raw ledgers are not archived); login-node micro-benchmarks of the released template; and four offline analyses of committed data:
\begin{itemize}[leftmargin=1.4em, itemsep=1pt]
\item \emph{Fault injection} (\RQ{1}): 19 mutation operators applied at every applicable site of the ten agent-written graphs, with each mutant run through the template's own self-test logic: \code{validate()} on the full and the smoke graph, placeholder filling for every cell, and the construction of every \code{sbatch} command. Twelve well-formedness operators (W1--W12) model slips that the checks are designed to catch; seven semantic operators (S1--S7) produce well-formed graphs with a different meaning.
\item \emph{Guard tests} (\RQ{1}): eleven tests of \I{1}--\I{4} and of the contracts in a sandbox copy of the template, with \code{nvidia-smi} and \code{scontrol} replaced by stubs, complemented by a requeue and guard experiment in Slurm jobs on \code{1g.10gb} MIG slices (Section~\ref{sec:rq1}).
\item \emph{Critical paths and waits} (\RQ{3}): the longest path of job run times through the pilot's dependency structure, and each job's wait, split into the dependency hold (from submission until its upstream jobs ended) and the scheduler wait (from then until the job started).
\item \emph{Check scaling} (\RQ{1}, \RQ{4}): the self-test logic applied to synthetic graphs of 1--32 chained stages with 1--1{,}024 cells per stage, taking the median of three repetitions.
\end{itemize}
The analyses use only the Python standard library and run in seconds on a login node. All statistics are descriptive: with one pilot and ten artifacts, we do not test hypotheses.

% =========================================================================
\section{Results, Part A: the checking layer}\label{sec:resA}

\subsection{RQ1: Static checking}\label{sec:rq1}

\begin{figure}[t]
\centering
\includegraphics[width=\textwidth]{figures/fig_mutation.pdf}
\caption{\RQ{1}. Faults injected into the ten agent-written graphs (columns R1--R10) and detected by the static checks, as detected/injected per graph; ``--'': operator not applicable. Rows W1--W12 are well-formedness faults (with the check designed to catch them); rows S1--S7 are well-formed graphs with a different meaning. Source: \code{generated/mutation\_results.csv}.}
\label{fig:mutation}
\end{figure}

\paragraph{Well-formedness faults} Figure~\ref{fig:mutation} shows the 641 mutants by operator and graph. The checks detected 438 of the 439 mutants of W1--W11: duplicate or unsafe stage names, dependencies on later or unknown stages, misspelt dependency kinds, missing GPU requests, misspelt placeholders in commands and contracts, removed variables, \emph{corr} edges between stages of different sizes, and steps without commands. Most of these faults were reported as validator diagnostics; placeholder faults were reported when the self-test filled a cell's commands. In one case, a duplicate stage name made the self-test abort with a traceback instead of a diagnostic, which still stops the submission but explains less. Checking one graph took 0.5--1.2\,ms.

\paragraph{Two defects} The study exposed two defects of the template. First, \emph{duplicate step names} within a stage (W12) went undetected in all 65 mutants, although the schema documentation requires step names to be unique. The consequence is worse than a missing diagnostic. Completion markers are keyed by step name, so the second step with a given name finds the first one's marker and is \emph{silently skipped}, already on the first execution and not only after a restart. Second, the one W9 mutant that escaped detection exposed a \emph{shadowed variable}. The RetriGen graph defines \code{WORK} (Listing~\ref{lst:graph}), and so does the site file \code{cluster.env}; when the mutant removes the graph's variable, \code{\{WORK\}} silently resolves to the site's directory instead of failing. %Both defects are cheap to fix, with a uniqueness check for step names and a warning when a graph placeholder resolves from the site file (Section~\ref{sec:impl-extend}); neither fix has been applied to the released skill yet.

\paragraph{Semantic faults} None of the 137 semantic mutants was detected (rows S1--S7). These mutants weaken a gate from \emph{ok} to \emph{any}; remove a dependency, a \emph{critical} flag or a data contract; swap two placeholders; remove the smoke overrides; or declare a GPU stage CPU-only. Each resulting graph is well-formed but means something other than the original. This is the boundary of static checking as \tool{} implements it. Some of these faults would surface later: a removed dependency, for example, causes a missing-input failure if the downstream step declares its inputs. Others, such as a weakened gate or swapped placeholders, would run without any signal. Such faults are the reason for the smoke run and for human review of the graph.

\paragraph{Guards} All eleven guard tests behaved as specified (\code{generated/guard\_tests.csv}). The runner refused to start outside a job (code 10, \I{1}). A step whose input pattern matched nothing failed with code 6, and a step that exited with 0 without producing its declared output failed with code 7. A rerun skipped a finished step and appended no second record (\I{3}). The GPU check refused to proceed when no GPU was visible (\I{2}). On the pre-limit signal, the handler requeued its own array task, or exited with code 99 once the requeue cap was reached (\I{4}). Two tests document limits rather than defects: a GPU that is visible but never used passes \I{2}, and a single file satisfies a pattern meant to match many. The sandbox test of \I{4} exercises only the handler's logic, with a stubbed \code{scontrol}; we therefore repeated it on Slurm.

\paragraph{Requeue and guards on compute nodes} We then tested \I{4} on Slurm with a dedicated graph (\code{docs/ase/analysis/i4test}). Its \code{requeue} stage appends 6{,}000 units of 0.05\,s of busy computation each (no \code{sleep}) on a \code{1g.10gb} slice with a 19-minute time limit, so that Slurm's pre-limit signal arrives about four minutes after the start. In the first run, the handler fired after 213\,s, wrote ``requeued'' to the job-ending ledger and called \code{scontrol requeue}, which failed inside the job with ``Incompatible Slurm plugin \ldots\ version (25.11.9)''. The Slurm clients on the inherited \code{PATH} did not match the server on that node. A diagnostic job confirmed this site-level version skew: every client on the \code{PATH} and in \code{/opt/software/slurm/current/bin} failed, and only the one in \code{/opt/software/slurm/25.11.8/bin} worked. Because the handler does not check the exit status of \code{scontrol}, the job exited with code 0, and Slurm recorded it as COMPLETED, with 4{,}029 of 6{,}000 units done and no restart. A dependent stage attached by an \emph{ok} edge would have started on incomplete data. This is a third defect of the template, independent of the skew: a failed requeue must fail the job.

In a second run, we prepended a working client to \code{PATH} in \code{cluster.env}, a change of configuration only. The handler then requeued the job after 223\,s (state REQUEUED, same job ID). The second attempt skipped the 4{,}247 recorded units and wrote the remaining 1{,}753, and the verifier found all 6{,}000 units exactly once, with no duplicates and no invalid digests, as Property~\ref{prop:resume} predicts. The first run's guard stage also confirmed two sandbox results on a compute node. With \code{CUDA\_VISIBLE\_DEVICES} emptied, the GPU check refused to proceed, and with the MIG slice visible it passed. In a stage with two steps named \code{dup}, the second step was skipped, the output file holding only the first step's line.

\paragraph{Cost} The time of the checks grows linearly with the number of cells. A least-squares fit over the 36 graph sizes gives 3.7\,ms plus 50\,$\mu$s per cell ($R^2=0.999$), that is, 1.6\,s for 32 stages of 1{,}024 cells each (Figure~\ref{fig:overhead}a).

\begin{quote}\small\textbf{Answer to \RQ{1}.} The static checks caught 438 of 439 well-formedness faults of eleven kinds, in under 2\,ms per real graph. They caught no duplicate step names, which can silently skip a step, and none of 137 well-formed but wrong graphs. On Slurm, requeue and resumption lost no work once a site-level client mismatch was worked around; without the workaround, the handler reported success although the job was unfinished.\end{quote}

\subsection{RQ2: Fault containment}\label{sec:rq2}

\begin{table}[t]
\centering\footnotesize
\caption{\RQ{2}, \RQ{5}. The 14 faults observed in the pilot and the artifact study, by root cause and by the layer that reported them. ``Cheapest layer'' is the earliest layer that could have caught the fault, in hindsight. Times are run times of the reporting job; H100-h are H100-equivalent hours.}\label{tab:faults}
\setlength{\tabcolsep}{3pt}
\begin{tabularx}{\textwidth}{@{}l >{\raggedright\arraybackslash}X l r r >{\raggedright\arraybackslash}p{2.1cm}@{}}
\toprule
Source & Cause & Reported by & Time & H100-h & Cheapest layer \\
\midrule
template & placeholders inside variables not expanded (T1) & self-test & 0.5\,s & 0 & self-test \\
template & \code{PIP\_FIND\_LINKS} hid the wheelhouse (T2; LLM4Perf) & gate & 11\,s & 0.0004 & login-node wheel check \\
environment & wheel unavailable (ProvenanceGuard, RetriGen, ExeCRE, BinarySPD) & gate & 12--22\,s & $\le$0.0009 & login-node wheel check \\
artifact & dataset loader asks an interactive question (ExeCRE) & gate & 224\,s & 0.0089 & gate \\
agent & comparison script expects a missing column (RetriGen) & gate & 1{,}099\,s & 0.0436 & gate \\
agent & pipeline exits with SIGPIPE, code 141 (LLM4Perf) & gate & 2{,}666\,s & 0.1058 & gate \\
pilot & model download broke off (3 submissions) & gate & 55--983\,s & 0.015--0.273 & gate \\
pilot & two jobs at 2.3\% and 5.0\% GPU utilization & GPU samples & -- & -- & smoke graph with every model \\
\bottomrule
\end{tabularx}
\end{table}

Table~\ref{tab:faults} lists all 14 faults we observed: 3 in the pilot's setup gates, 8 in the artifact study's gates, 1 in a self-test, and 2 in GPU telemetry. Four of them originated in the tooling under study: two template defects and two errors in the agent's configuration. The others came from the environment, the artifacts and the pilot's own infrastructure.

\paragraph{Gates} All 11 failing gates stopped their submissions before any main stage started. Slurm cancelled every dependent attached by \emph{ok} or \emph{corr} edges before it was allocated: 75 jobs in the three failed pilot submissions (25 each: 14 matrix, 6 side-study, 3 repair, 1 coding and 1 analysis job) and 12 in the artifact study, 87 in all. None of these jobs used a GPU-second. The failing gates themselves cost 0.015--0.273 GPU-hours each in the pilot, where they ran on whole H100s, and 0.0004--0.106 H100-equivalent hours each in the artifact study, where they ran on MIG slices. Two further pilot submissions were cancelled by the operator during development after their setup jobs had run for 256\,s and 872\,s on whole H100s at 0\% GPU utilization; we do not count them as gate events.

Collectors behaved differently, by design. Because they are attached with \emph{any} edges, the 8 collectors behind failed artifact gates were not cancelled; 7 ran, mostly failing on the same environment problem as their gate, for 0.0078 H100-equivalent hours together. How much a gate saves depends on what its dependents would have done without it. The pilot's dependents requested 29 H100s with time limits of up to 24 hours, an upper bound of 696 GPU-hours per submission. This bound comes from the pilot configuration, not from the archived accounting data, and it overstates the realistic saving, because most dependents would likely have failed early on the missing models. %We did not measure the realistic saving.

\begin{figure}[t]
\centering
\includegraphics[width=\textwidth]{figures/fig_faults.pdf}
\caption{\RQ{2}. (a) Time to feedback of the 11 failed setup gates: the gate's run time, and the time from submission to failure including the queue wait. (b) GPU cost of each failed gate against its run time, by root cause; dashed and dotted lines: cost of a whole H100 and of a \code{1g.10gb} slice. Sources: \code{generated/gate\_feedback.csv}, \code{generated/faults.csv}.}
\label{fig:faults}
\end{figure}

\paragraph{Time to feedback} Figure~\ref{fig:faults}a separates two durations that are easily conflated: how long a failing gate runs, and how long the researcher waits for its verdict. The failing artifact gates ran for a median of 22\,s, but the researcher learned of the failure a median of 10.4 minutes after submission (4.4--55.8 minutes), because each gate first waited in the queue (median 10.1 minutes). For half of these failures, the wait bought nothing. The four unavailable wheels, and plausibly the template defect T2, could have been caught in seconds on the login node by resolving the requirements against the wheelhouse without installing anything. We added such a check (\code{pip install -{}-dry-run -{}-no-index}) to the evaluation harness \emph{during} the study, after these failures; it is not part of the evaluated template.

\paragraph{Cost per fault} Figure~\ref{fig:faults}b shows that the cost of a failed gate is its run time multiplied by its GPU share. Running the artifact gates on \code{1g.10gb} slices rather than on whole H100s, as the pilot did, cuts the cost of a failure of equal duration sevenfold. A gate that only installs packages and checks files does not need a whole GPU.

\begin{quote}\small\textbf{Answer to \RQ{2}.} Setup gates stopped all 11 failing submissions before any main stage started, and Slurm cancelled all 87 dependents attached by \emph{ok} or \emph{corr} edges at zero GPU cost, while collectors attached by \emph{any} edges still ran. Feedback took about ten minutes, almost all of it queue wait, and four or five of the eight artifact-gate failures could have been caught on the login node.\end{quote}

\subsection{RQ3: Concurrency and observability}\label{sec:rq3}

\begin{table}[t]
\centering\footnotesize
\setlength{\tabcolsep}{4pt}
\caption{\RQ{3}. Smoke runs of the pilot (10 jobs and 36 steps each, all completed). Concurrency is summed job run time over makespan; the bound is summed job run time over the critical path of job run times; waits are medians over non-gate jobs, split into dependency hold and scheduler wait.}\label{tab:smoke}
\begin{tabular}{@{}l r r r r r r r@{}}
\toprule
Run & Nodes & Makespan (s) & Crit.\ path (s) & Concurrency & Bound & Ratio & GPU-h \\
\midrule
1 & 7 & 1{,}068 & 1{,}021 & 2.68 & 2.81 & 0.96 & 1.12 \\
2 & 8 & 1{,}461 & 1{,}155 & 2.05 & 2.60 & 0.79 & 1.19 \\
3 & 7 & 1{,}113 & 1{,}070 & 2.39 & 2.48 & 0.96 & 1.06 \\
\bottomrule
\end{tabular}
\end{table}

\begin{figure}[t]
\centering
\includegraphics[width=\textwidth]{figures/fig_timeline.pdf}
\caption{\RQ{3}. Jobs running over time, from Slurm accounting. (a) The three smoke runs of the pilot, with mean concurrency and the critical-path bound. (b) The full pilot run, jobs running and pending, until the snapshot 3.1 hours after its first job started. Source: \code{generated/timeline.csv}.}
\label{fig:timeline}
\end{figure}

\paragraph{Concurrency} \tool{} queues all stages at once and lets Slurm resolve their dependencies, so independent jobs run side by side and their queue waits overlap. Table~\ref{tab:smoke} and Figure~\ref{fig:timeline}a show the pilot's three smoke runs. Each completed all 10 jobs and 36 steps on 7--8 nodes in 18--24 minutes, against 44--50 minutes of summed job run time, which gives a mean concurrency of 2.05--2.68. A speed-up over serial execution would be an uninformative comparison here, because any submission of the same jobs with Slurm dependencies would also run them concurrently. The informative comparison is with the critical path, the fastest that any execution of this graph could finish if every job started the moment it became eligible: 1{,}021--1{,}155\,s (17--19 minutes). The smoke runs reached 79--96\% of that bound. Their non-gate jobs waited a median of 676\,s for their upstream jobs and only 82\,s in the scheduler queue. Under the load we observed, the graph's shape, not the queue, therefore determined the makespan.

The full run (Figure~\ref{fig:timeline}b) is a snapshot taken 11{,}140\,s (3.1 hours) after its first job started, while the run was still in progress; its final totals are not part of this article. By then, 19 of 26 jobs had started on 8 nodes, with up to 18 jobs (21 GPUs) running at once and a mean concurrency of 14.6 jobs; 18 steps had completed and none had failed; and the jobs had used 51.8 GPU-hours. The non-gate jobs had waited a median of 1{,}774\,s, of which 845\,s was dependency hold and 928\,s scheduler wait.

\paragraph{Observability} In the scheduler's view (\code{squeue} and the default \code{sacct} fields), all 14 matrix jobs of the full run were running normally. The 60-second GPU samples recorded alongside \I{2} tell a different story (Figure~\ref{fig:gpu}, 159--176 samples per job). Twelve jobs averaged 73.3\% utilization (46.8--85.0\%) and were above 10\% utilization in 81--99\% of their samples. The two DeepSeek-Coder-V2 jobs averaged 5.0\% and 2.3\%, were above 10\% in only 22\% and 10\% of their samples, and had been in that state for 2.7 hours at the time of the snapshot. Because they were not failing, no other signal flagged them. We did not diagnose the cause; we note only that DeepSeek-Coder-V2, unlike the other models, is a mixture-of-experts model. The smoke graph could not have revealed this anomaly, because it ran only the first model.

\begin{figure}[t]
\centering
\includegraphics[width=\textwidth]{figures/fig_gpu.pdf}
\caption{\RQ{3}. GPU telemetry of the 14 matrix jobs of the full pilot run, by model (ordered by total parameters) and benchmark: (a) mean utilization, (b) share of the 60-s samples above 10\% utilization. The outlined row is the anomaly that scheduler states did not show. Source: \code{data/pilot/gpu\_util.csv}.}
\label{fig:gpu}
\end{figure}

Among the other six models, utilization tended to be higher for larger models, rising from 46.8--50.5\% at 7B parameters to 84.0--85.0\% at 32B. We report this trend descriptively only. A rank correlation over the six model means ($\rho=0.83$) is fragile for three reasons: it rests on excluding the outlier after seeing it (with DeepSeek-Coder-V2 included, $\rho=0.64$); it ranges from 0.70 to 1.00 when any one model is left out; and model size is confounded with model family and serving configuration. One untested explanation is that small models leave the GPU waiting while the harness runs tests on the CPU.

\begin{quote}\small\textbf{Answer to \RQ{3}.} The smoke graphs reached 79--96\% of the concurrency their dependency structure permits. This result reflects Slurm and the graph's shape; we did not compare it with an alternative. Per-job GPU samples exposed two jobs that the scheduler reported as running but that idled at 2--5\% utilization for 2.7 hours; the cause was not diagnosed.\end{quote}

\subsection{RQ4: Overhead}\label{sec:rq4}

\begin{table}[t]
\centering\footnotesize
\caption{\RQ{4}. Login-node cost of the template (\code{data/pilot/microbench.csv}, \code{generated/microbench\_checks.csv}). $n$: repetitions per value.}\label{tab:micro}
\begin{tabular}{@{}>{\raggedright\arraybackslash}p{4.6cm} >{\raggedright\arraybackslash}p{3.1cm} r r@{}}
\toprule
Operation & Size & Time & $n$ \\
\midrule
Scaffold \code{hpc/} into a repository & 18 files & 0.18\,s (median) & 3 \\
Self-test (\V{1}--\V{6}, full and smoke graph, dry runs) & 4 stages, 10 + 6 cells & 0.48\,s (median) & 3 \\
Static checks only (synthetic graphs) & 1--32{,}768 cells & 3.7\,ms + 50\,$\mu$s/cell & 3 \\
Runner overhead per executed / skipped step & 50 steps & 7.1 / 1.2\,ms & 1 \\
Audit snapshot & 10 / 100 / 1{,}000 / 5{,}000 cells & 0.02 / 0.03 / 0.12 / 0.52\,s & 1 \\
Collect (tables, summaries, Markdown and \LaTeX{}) & 10{,}000 / 100{,}000 records & 0.17 / 1.41\,s & 1 \\
\bottomrule
\end{tabular}
\end{table}

\begin{figure}[t]
\centering
\includegraphics[width=\textwidth]{figures/fig_overhead.pdf}
\caption{\RQ{1}, \RQ{4}. (a) Time of the static checks over a grid of graph sizes (median of 3 per point; contour lines interpolate between the 36 measured points). (b) Login-node cost of the audit, collector and static checks against workload; dashed: the runner's per-step overhead. Sources as in Table~\ref{tab:micro}.}
\label{fig:overhead}
\end{figure}

Table~\ref{tab:micro} and Figure~\ref{fig:overhead} show what the template's machinery costs. Scaffolding takes 0.18\,s and the self-test 0.48\,s. The runner adds 7.1\,ms per executed step and 1.2\,ms per skipped one, which is negligible for GPU jobs that run for minutes to hours. An audit snapshot costs about 17\,ms plus 0.1\,ms per cell (a least-squares fit over four sizes with single measurements), or 0.52\,s at 5{,}000 cells, and collecting 100{,}000 records takes 1.4\,s. These costs are small enough for the login node, where the audit and the checks run. The substantial cost is the smoke run itself: 1.06--1.19 GPU-hours per pilot smoke run.

\begin{quote}\small\textbf{Answer to \RQ{4}.} The machinery costs milliseconds per step and at most seconds per check, audit or collection. The dominant cost is the smoke run, about one GPU-hour per pilot smoke run.\end{quote}

% =========================================================================
\Needspace{12\baselineskip}
\section{Results, Part B: the agent on unseen artifacts}\label{sec:resB}

\subsection{RQ5: Feasibility}\label{sec:rq5}

\begin{table}[t]
\centering\footnotesize
\caption{\RQ{5}. The ten artifacts: configuration written by the agent session, data contracts declared, setup gates in submission order (\xmark{} failed, \cmark{} passed, -- withdrawn before start) and H100-equivalent hours (summed from unrounded job values). LoC: non-generated lines of the graph, comparison script and hooks; glue: additional scripts between documented steps.}\label{tab:artifacts}
\setlength{\tabcolsep}{3pt}
\begin{tabular}{@{}l l l r r r c l r@{}}
\toprule
Artifact & Venue & Mode & Stages/steps & LoC & Glue & Contracts & Gates & H100-h \\
\midrule
SemGuard~\cite{wang2025semguard} & ASE'25 & probe & 3/12 & 149 & -- & -- & -- & -- \\
ProvenanceGuard~\cite{she2026provenanceguard} & ASE'26 & subset & 3/6 & 162 & -- & -- & \xmark\cmark & 0.0074 \\
RetriGen~\cite{zhang2025retrigen} & TOSEM'25 & subset & 4/14 & 130 & 42 & N, P & \xmark\xmark & 0.0465 \\
ExeCRE~\cite{dong2026execre} & ASE'26 & subset & 4/9 & 225 & 18 & P & \xmark\xmark\cmark & 0.0239 \\
EffiSkel~\cite{yu2026effiskel} & FSE'26 & probe & 3/15 & 161 & -- & -- & -- & -- \\
BinarySPD~\cite{li2026binaryspd} & TOSEM'26 & subset & 3/7 & 164 & -- & P & \xmark\cmark & 0.0044 \\
IRCoCo~\cite{li2024ircoco} & FSE'24 & probe & 3/12 & 149 & -- & -- & -- & -- \\
RobustExt~\cite{rabbi2026robustness} & EMSE'26 & subset & 3/8 & 116 & 45 & P & \cmark & 0.0042 \\
LLM4Perf~\cite{wang2026llm4perf} & ICSE'26 & subset & 3/5 & 54 & -- & P & \xmark\xmark & 0.1072 \\
LogReasoner~\cite{ma2026logreasoner} & TOSEM'26 & probe & 3/9 & 137 & -- & -- & -- & -- \\
\midrule
Total & \multicolumn{2}{l}{6 subset, 4 probe} & 32/97 & 1{,}447 & 105 & 5 of 10 & 8\xmark\,4\cmark & 0.1936 \\
\bottomrule
\end{tabular}
\end{table}

\paragraph{Configuration} Table~\ref{tab:artifacts} summarises the ten configurations. They have 3--4 stages and 5--15 steps each, 32 stages and 97 steps in total, and take 54--225 lines each, counting the graph, the comparison script and the model-server hooks (1{,}447 lines in all, most of them in the graphs). All ten passed the self-test before submission. Three needed 18--45 lines of \emph{glue} code between documented steps. Although no artifact file was edited, this count understates the adaptation: the RetriGen glue also works around an indexing bug in the artifact's own script, and several environments deviate from the published ones (Table~\ref{tab:interventions}).

The agent used the graph's expressive features sparingly. Every subset graph has a setup gate (\emph{ok}) and a collector (\emph{any}); only RetriGen uses \emph{corr} and declares inputs; and outputs are declared by 16 of 97 steps in five graphs. The four probe graphs attach their probe stage with an \emph{any} edge. This choice is deliberate, so that every documented entry point runs even when setup fails, but it also means that their setup stage gates nothing.

\paragraph{Gates} The six subset configurations were smoke-tested; Table~\ref{tab:faults} and Section~\ref{sec:rq2} report their faults. Including resubmissions, 12 setup gates ran and 8 failed: four on unavailable wheels, one on the template defect T2, one on an artifact's interactive dataset prompt, and two on the agent's own configuration (a pipeline that exited with SIGPIPE, and a comparison script that expected a column the data lacked). Fixes were resubmitted after a median of 8.3 minutes (2.7--13.3, $n=6$). After this round, four of the six subset artifacts had passed their gate, at a total cost of 0.19 H100-equivalent hours. All queued jobs, including the four probes, were then withdrawn before any main stage or probe step started. Three days later, we resubmitted the four passing configurations in smoke mode, with the template unchanged: all four setup gates passed again, in 41--214\,s. Their main-stage smoke jobs, on \code{3g.40gb} slices and one whole H100, were queued but had not started when this article was prepared, so we report no main-stage outcome.

\paragraph{Defects found while configuring} While reading the artifacts to configure them, the sessions reported suspected defects in all four probe artifacts, such as a missing \code{requirements.txt}, a hard-coded placeholder path, or an import of a module that does not exist. They also reported missing steps in three subset artifacts, for example the absence of code that converts the retrieval outputs into the files the evaluation script reads. LLM4Perf ships no evaluation code, so its comparison step reports the paper's numbers as \emph{not comparable} instead of passing silently. These are suspicions, not confirmed defects: the probes that would confirm or refute them by exit code have not run.

\begin{table}[t]
\centering\small
\caption{\RQ{5}. Human interventions and deviations from the published artifacts (from \code{data/study\_log.md}). Actor: H = human (author), A = agent session, T = template fix.}\label{tab:interventions}
\begin{tabularx}{\textwidth}{@{}l >{\raggedright\arraybackslash}X l@{}}
\toprule
Artifact(s) & Intervention or deviation & Actor \\
\midrule
all & executed subset and comparison target chosen per artifact & H \\
all & extra wheels downloaded from PyPI into a study wheel directory where the wheelhouse lacked them (e.g.\ \code{crowd-kit}, \code{javalang}, exact pins of LLM4Perf) & H \\
vLLM-based (3) & \code{opencv} module added to satisfy the wheelhouse's vLLM & H \\
RetriGen & \code{numpy} 1.26.2 $\to$ 1.26.4 (no wheel); \code{torch} 1.13.1 $\to$ 2.1.2; \code{codebleu} package instead of the authors' script; glue converts intermediate files and works around an off-by-one in \code{embed.sh} & H, A \\
RobustExt & \code{transformers} 4.21.1 $\to$ 4.45.2 (to load Qwen2.5); hard-coded model paths mapped in glue; metrics functions loaded from a module-level script & A \\
SemGuard, EffiSkel, IRCoCo & Python 3.10 instead of 3.11 (a test-harness dependency does not build on 3.11), IRCoCo instead of the required 3.7; unused \code{datasets} dependency dropped & H \\
ExeCRE & \code{HF\_DATASETS\_TRUST\_REMOTE\_CODE=1} set for the dataset loader (no code change) & H \\
LLM4Perf & evaluation code absent: comparison marked not comparable & A \\
all & template defects T1, T2 fixed during the study; login-node wheel check added to the harness & T, H \\
\bottomrule
\end{tabularx}
\end{table}

\paragraph{Human effort} Table~\ref{tab:interventions} lists every intervention recorded in the study log. Its length is itself a finding. Bringing ten published artifacts to the point of a passing setup stage on an academic cluster required environment substitutions for five of them (Python versions for three, library versions for two), one of which changes the major version of a deep-learning library (\code{torch} 1.13 to 2.1). Such substitutions could change the published numbers and must be reported with any reproduction claim. We did not record the agent sessions' transcripts, duration or token use, or the number of times the human corrected a session; the log covers only interventions that changed the environment or the configuration.

\begin{quote}\small\textbf{Answer to \RQ{5}.} With the skill, a new agent session per artifact produced graphs of 54--225 lines that passed the static checks for all ten artifacts. Four of the six subset artifacts passed their setup gate, one at the first attempt and three after one or two fixes, at a cost of 0.19 H100-equivalent hours and with the human interventions of Table~\ref{tab:interventions}. Whether the main stages run and reproduce the papers' numbers is not yet known.\end{quote}

% =========================================================================
\section{Discussion}\label{sec:discussion}

\paragraph{What the checks are for} The fault-injection study draws a clear line. Static checks are cheap and reliable for the faults they target: the slips that an agent, or a person, makes when writing a graph by hand, such as a misspelt placeholder, a stage listed out of order or a missing GPU request. None of the 13 faults that reached the cluster was of this kind. This is expected, since every graph had passed the self-test. Nor did any of them belong to the two well-formedness classes that escaped the checks, duplicate step names and shadowed variables. The faults that did reach the cluster were environmental (wheels, downloads, an interactive prompt) or semantic (a script expecting a missing column), and they were caught by executing a cheap stage on real nodes. The three defects we found, a silently skipped duplicate step, a shadowed variable and a requeue handler that reports success when the requeue fails, carry a lesson for tool builders: a checker needs its own fault-injection tests, and run-time handlers need tests on the real scheduler. The requeue handler passed its stub test and failed in its first real run, and the schema documentation required unique step names that no code enforced.

\paragraph{Where gates belong} The gates did what Slurm's dependency semantics promise, but the queue made them slow: about ten minutes of waiting for 22 seconds of work. Half of the gate failures were package-resolution failures that a login-node check finds in seconds without a GPU. Checks should therefore be ordered by their cost and by where they can run: dependencies resolved on the login node; file, model and pre-flight checks in a gate on the smallest GPU slice; and whole GPUs reserved for the main matrix.

\paragraph{What smoke runs cover} A smoke run that executes every stage, step and edge once is a meaningful test of the \emph{pipeline}, and in our data it exposed every fault that a gate could expose. It is not a test of the \emph{experiment's input space}: the first $k$ cells of the pilot covered one of seven models, and the one model-specific anomaly lay outside them. A projection that covers every value of every matrix parameter at least once would have included DeepSeek-Coder-V2, at roughly seven times the smoke cost. Whether it would have exposed the anomaly within a smoke run's few minutes is untested.

\paragraph{Telemetry as part of the record} The two idle GPU jobs would not have changed any number in the pilot study, but they did add to its cost, and a reader of the study would want to know about them. Recording utilization next to the results is cheap. As implemented, \I{2} observes utilization rather than enforcing it; an in-job threshold that flags or stops a job whose utilization stays below $x$\% for $T$ minutes is a natural extension.

\paragraph{When not to use \tool{}} \tool{} targets Slurm clusters, GPU-bound experiments with a few stages and parameter grids, and users who want plain Slurm jobs that cluster staff can read. Workflows with many fine-grained tasks, complex data dependencies, or a need to run on several schedulers or clouds are better served by a workflow system, which a \tool{} step can still invoke. The template can also be configured by hand, without an agent; nothing in it depends on Claude Code.

% =========================================================================
\section{Related work}\label{sec:related}

\begin{table}[t]
\centering\small
\footnotesize
\caption{\tool{} and related systems. \cmark: provided; \pmark: partially or optionally; --: not provided. The comparison is of documented features, not an empirical evaluation.}\label{tab:related}
\setlength{\tabcolsep}{3pt}
\begin{tabular}{@{}>{\raggedright\arraybackslash}p{4.3cm} c c c c c c@{}}
\toprule
& Task & File & Re- & Slurm & Smoke & Agent- \\
& graph & contracts & sume & native & subset & authored \\
\midrule
Snakemake~\cite{molder2021snakemake}, Nextflow~\cite{ditommaso2017nextflow} & \cmark & \cmark & \cmark & \pmark & \pmark & -- \\
Pegasus~\cite{deelman2015pegasus}, Parsl~\cite{babuji2019parsl}, FireWorks~\cite{jain2015fireworks} & \cmark & \pmark & \cmark & \pmark & -- & -- \\
CWL~\cite{crusoe2022cwl} & \cmark & \cmark & \pmark & \pmark & -- & -- \\
signac~\cite{adorf2018signac} & \pmark & \pmark & \cmark & \pmark & -- & -- \\
Sacred~\cite{greff2017sacred}, Sumatra~\cite{davison2012sumatra} & -- & -- & -- & -- & -- & -- \\
ExecutionAgent~\cite{bouzenia2025executionagent}, Repo2Run~\cite{hu2025repo2run} & -- & -- & -- & -- & -- & \cmark \\
\tool{} & \cmark & \pmark & \cmark & \cmark & \cmark & \cmark \\
\bottomrule
\end{tabular}
\end{table}

\paragraph{Workflow and experiment management} Workflow systems such as Pegasus~\cite{deelman2015pegasus}, Snakemake~\cite{molder2021snakemake}, Nextflow~\cite{ditommaso2017nextflow}, Parsl~\cite{babuji2019parsl}, FireWorks~\cite{jain2015fireworks} and the Common Workflow Language~\cite{crusoe2022cwl} describe computations as graphs of tasks with file dependencies, can resume from completed tasks, and submit to schedulers including Slurm. signac~\cite{adorf2018signac} manages parameter spaces and job state on HPC. Sacred~\cite{greff2017sacred} and Sumatra~\cite{davison2012sumatra} record configuration and provenance of runs, and Sumatra links figures in a paper to the computations that produced them. Table~\ref{tab:related} compares their documented features. \tool{} borrows the central abstraction of these systems, a task graph with file contracts, and adds little to it as a workflow language. It differs in its target and in its division of labour: a small JSON schema meant to be written by an agent and checked by a fixed validator; compilation to plain Slurm arrays and dependencies with site guards (no login-node work, a GPU behind every job); a smoke graph derived from the full one; and an audit that joins scheduler and pipeline state. The design assumes that an agent writes correct graphs more easily in this schema than in Snakemake or Nextflow; this assumption is untested.

\paragraph{Reproducibility infrastructure} Environment capture with ReproZip~\cite{chirigati2016reprozip} or containers~\cite{kurtzer2017singularity}, the Popper convention~\cite{jimenez2017popper}, artifact badges~\cite{acm2020badges}, checklists~\cite{pineau2021improving} and data principles~\cite{wilkinson2016fair} make computations rerunnable and reviewable. Large studies document how often computations are not rerunnable~\cite{collberg2016repeatability,pimentel2019jupyter,trisovic2022research,gundersen2018state}, and SE studies show the same for research artifacts~\cite{timperley2021artifact,liu2024artifacts}. Many of the reported failures are mundane, which motivates cheap checks before expensive runs. Guidelines for LLM studies in SE~\cite{sallou2024breaking,baltes2025guidelines} and for reporting parallel performance~\cite{hoefler2015scientific} ask for complete reporting of what was run. Boa~\cite{dyer2013boa} made ultra-large-scale repository mining accessible to SE researchers; \tool{} has a similar aim for GPU experiments, an aim that we have not yet evaluated with users.

\paragraph{Agents that set up and reproduce research code} SUPER~\cite{bogin2024super}, CORE-Bench~\cite{siegel2024corebench}, PaperBench~\cite{starace2025paperbench} and MLAgentBench~\cite{huang2024mlagentbench} benchmark agents that set up repositories, reproduce computational results or run machine-learning experiments, mostly in a single container. ExecutionAgent~\cite{bouzenia2025executionagent}, Repo2Run~\cite{hu2025repo2run} and EnvBench~\cite{eliseeva2025envbench} target building environments and executing arbitrary projects, the step at which half of our gate failures occurred. Closest to our Part B, Artisan~\cite{baek2026artisan} has an agent reproduce the tables of SE papers from their artifacts by generating inspectable scripts, and it found errors in both papers and artifacts. These works measure what an agent \emph{can} do in a sandbox. \tool{} instead constrains and checks what an agent's configuration \emph{may} do on shared, scheduled infrastructure, and it does not yet measure reproduction success. Runtime enforcement of agent behaviour, as in AgentSpec~\cite{wang2026agentspec}, is complementary: it governs the agent's actions, whereas \tool{} checks the configuration the agent hands over. LLMs that write parallel HPC code~\cite{nichols2024hpccoder} address the programs, not their orchestration.

% =========================================================================
\section{Threats to validity}\label{sec:threats}

\paragraph{Construct validity} Mean concurrency and the critical-path bound describe the submitted graph and Slurm, not a benefit of \tool{} over another way of submitting the same jobs. The 696 GPU-hour bound of Section~\ref{sec:rq2} assumes that every dependent would have run to its time limit; the realistic saving is unknown. GPU utilization as reported by \code{nvidia-smi} measures the share of time in which a kernel was running, not useful throughput, and the 10\% ``busy'' threshold is arbitrary. The requeue experiment used a synthetic busy-computation workload on a MIG slice, one run per condition, and a client workaround for a site-level Slurm version skew; it shows that the mechanism can work, not how reliably it works. Counts of configuration lines depend on formatting. ``Time to feedback'' measures when a failure was recorded, not when a person noticed it.

\paragraph{Internal validity} The author of this article is the developer of \tool{}, wrote the pilot graph, performed the artifact selection as the only screener, decided the executed subsets, and classified the faults; no second rater checked these decisions. The pilot pipeline is the prototype of the template, not the released template. The treatment changed during the artifact study: two template defects were fixed and a login-node wheel check was added to the harness. The fault-injection operators were designed by us and may not reflect the distribution of real agent errors; detection rates are per operator and should not be averaged into one number. Each artifact was configured in a single agent session, so we cannot separate the skill's effect from session-to-session variation, and the sessions' transcripts, model version and token use were not recorded.

\paragraph{External validity} All runs used one cluster, one scheduler and one agent. Other Slurm sites need different GPU names, module stacks and wheel set-ups, and other schedulers need a different submitter and audit. The artifacts are LLM-related, Python-based SE artifacts that passed feasibility screening; the ranking's bonus for local model serving shaped which repositories were screened (14 of the 24 screened repositories use Ollama, against 22 of the 118 ranked), although only 3 of the 10 selected use it. Multi-node training jobs, which exercise \tool{}'s multi-node launch, were not among the subjects.

\paragraph{Conclusion validity} All results are descriptive, from one pilot, three smoke runs, one partial full run and ten artifacts. The full-run numbers are a snapshot taken while the run was in progress. The pilot's raw step ledgers and GPU samples are not archived; numbers derived from them (step counts, utilization means) are reproducible only from the committed summaries.

\paragraph{Deviation from the planned protocol} An earlier version of this work announced a multi-repository study with four domains (LLM evaluation, vision training, graph learning and SE), a no-skill baseline, a human rater of the claim-to-job tables, and full runs. That plan was not publicly registered. The study we report deviates from it in every one of these points: it covers only SE artifacts, has no baseline and no rater, and stops at the setup gates.
% TODO(author): if the earlier version is publicly available (e.g., a preprint), cite it here; if it is not, consider whether this paragraph is needed, since readers cannot consult the earlier plan.

% =========================================================================
\section{Conclusion and future work}\label{sec:concl}

Coding agents can generate experiment configurations, but generating a plausible configuration is not the same as establishing that an experiment is safe to execute, correctly orchestrated, or capable of reproducing published results. This distinction matters on shared high-performance computing (HPC) clusters, where configuration errors can consume scarce GPU resources, remain undetected behind scheduler queues, or produce misleading evidence of experimental success. This work addresses that gap through \tool{}, a validation and execution layer that translates agent-generated experiment plans into controlled Slurm workflows, checks their structural properties before submission, and uses inexpensive setup gates and smoke runs to detect failures before committing resources to full experiments.

Our evaluation provides evidence for the feasibility and practical value of this separation of responsibilities. In the graph-validation experiment, the checking layer detected 438 of 439 injected well-formedness faults before submission. The execution experiments further showed that setup gates can intercept failures arising from dependencies, environment configuration, and execution-time assumptions, while smoke runs can expose faults that static analysis cannot establish from the graph alone. The study also exposed two silently idle GPU jobs and demonstrated that scheduler-aware mechanisms can address execution problems that are otherwise difficult to diagnose from configuration files. Together, these findings support a central design principle: \textit{agent-generated experiment plans should be treated as proposals to validate, not as instructions to execute without independent checks.}

The results also identify important limits. A structurally valid graph can still encode an incorrect experiment, and static analysis cannot establish that its outputs support the claims being evaluated. The tool itself contained three defects, including a requeue-handling defect that its stub test did not reveal but a real-scheduler test exposed. Moreover, queue delays dominated the time to feedback, showing that inexpensive checks alone do not guarantee a responsive workflow on a shared cluster. Finally, although an agent configured ten software-engineering research artifacts and four of six subset configurations eventually passed their setup gates, the main experimental stages had not completed when this study was prepared. We therefore make no claim that \tool{} reproduces the published results of these artifacts. The evidence establishes a feasibility result for configuration, validation, and setup, not end-to-end experimental reproducibility or superiority over an alternative workflow.

Three next steps are necessary to establish those stronger claims. First, we will execute the main stages of the artifacts that passed setup validation and compare their outputs with the published results using comparison criteria and tolerances defined in advance. Second, we will compare agent sessions using \tool{} with sessions operating without it, repeating both conditions across multiple sessions to distinguish the effect of the tool from variation in agent behaviour. Third, we will independently assess the accuracy and completeness of the generated claim inventories and claim-to-job mappings, including whether each reported result is supported by an executed step and an available output. In parallel, we will address the three identified implementation defects, strengthen dependency checks on the login node, and extend fault-injection tests to cover the scheduler behaviours on which safe execution depends.

%Ultimately, \tool{} is not intended to make coding agents infallible or to replace scientific judgment. Its purpose is to make agent-assisted experimentation more inspectable, resource-conscious, and accountable by placing explicit checks between generated plans and shared infrastructure. Establishing whether this approach improves reproducibility, reduces human effort, and prevents wasted GPU time requires broader comparative studies and completed end-to-end replications. The present work provides the foundation for those studies and makes the boundary of the current evidence explicit.

% \tool{} lets a coding agent propose an experiment graph for a shared Slurm cluster and leaves every decision about what runs to fixed programs. Evaluated on its own, the checking layer caught 438 of 439 well-formedness faults in agent-written graphs before submission, stopped every failing submission in a cheap setup gate before its main stages, ran smoke graphs close to the concurrency their structure permits, and exposed two silently idle GPU jobs. Its overhead on the login node is milliseconds to seconds, and a smoke run costs about one GPU-hour. The same evaluation showed the layer's limits: no static check catches a well-formed but wrong graph; the tool itself had three defects, one of which only a test on the real scheduler exposed; and queue waits, not gate run times, dominate the time to feedback. With the skill, an agent configured ten SE artifacts up to passing or failing setup gates, with human help that we have documented.

%Three experiments would turn this feasibility evidence into evidence of benefit: (1) running the planned subsets of the four artifacts whose gates passed and comparing the outputs with the published numbers under tolerances fixed in advance; (2) a baseline in which the same agent configures the same artifacts without the skill, with several sessions per condition; and (3) an evaluation of the agent's claim inventories and claim-to-job tables against independent raters, together with a machine check that every row names a step whose declared output exists. On the tool side, we will fix the three defects (reject duplicate step names, warn when graph variables shadow site variables, and fail the job when \code{scontrol requeue} fails), add a login-node dependency check to the self-test, and report the Slurm client skew we encountered to the site's support team. %Section~\ref{sec:impl-extend} describes further extensions, from static contract linkage and template integrity checks to other schedulers and the needs of large SE testing campaigns.

% =========================================================================
\backmatter

\bmhead{Data availability}
The evaluation repository contains the selection data (all 4{,}814 publications, every automatic and manual decision), the ten artifact configurations, the Slurm accounting extracts, the pilot's derived CSV files, the scripts that regenerate the article's figures and derived numbers (\code{docs/ase/analysis/}), and the requeue and guard experiment (\code{docs/ase/analysis/i4test/}) with its accounting and evidence (\code{docs/ase/generated/reruns\_jobs.csv}, \code{i4\_verify.json}); \code{make check} verifies offline that the committed selection and study tables regenerate byte for byte. The pilot's raw ledgers and the agent sessions' transcripts are not included.
% TODO(author): archive the evaluation repository and the evaluated skill version on Zenodo and insert both DOIs.

\bmhead{Code availability}
\tool{} is open source at \url{https://github.com/ai4se4ai-lab/use-hpc-claude-skill}~\cite{usehpc2026repo}; the evaluated version is release 1.0.1 (commit \code{2fc1509}). The evaluation repository~\cite{usehpc2026eval}, available at \url{https://github.com/ai4se4ai-lab/use-hpc-evaluation}, contains the selection data (all 4{,}814 publications, every automatic and manual decision), the ten artifact configurations.

% TODO(author): the repository has no LICENSE file yet; add one and state the licence here.

\section*{Declarations}


\bmhead{Funding}
This research received no specific grant from any funding agency in the public, commercial, or not-for-profit sectors.

\bmhead{Ethical Approval}
Ethical approval was not required for this study as it involved only the analysis of publicly available open-source software repositories.

\bmhead{Informed Consent}
Informed consent was not applicable.

\bmhead{Competing interests} The author is the developer of \tool{}. The author declares no other competing interests.


\bmhead{Author contributions} Majid Babaei: Conceptualization, Methodology, Software, Validation, Formal analysis, Investigation, Data curation, Writing – original draft, Writing – review \& editing, Visualization. 

\bmhead{Use of generative AI} The approach studied here uses an LLM agent (Claude Code) to configure experiment pipelines, as described in Sections~\ref{sec:approach} and \ref{sec:subjects}.


\bibliography{bibliography}

\end{document}
